\subsection{Diamonds}
\label{sec:diamonds}

The aim of this section is to prove the following.

\begin{theorem}
\label{thm:diamond}
For any $\ell \in [4]$, {\sc Planar $\oldC_\ell$-Steiner Network} is $\mathsf{W}[1]$-hard parameterized by the number $k$ of terminals and 
does not admit an $f(k) \cdot n^{o(\sqrt{k})}$ algorithm for any computable function $f$,
unless $\mathsf{ETH}$ fails.
\end{theorem}

We only formally prove the statement for pure out-diamonds as it will become clear from the proof that
to handle 
\begin{enumerate}
\item flawed out-diamonds, it suffices to add a vertex $s$ and edges $(s,r_1)$ and $(s,r_2)$ in the construction below;
\item pure in-diamonds, it suffices to reverse the direction of every edge in the construction below; and
\item flawed in-diamonds, it suffices to add a vertex $t$ and edges $(r_1,t)$ and $(r_2,t)$, and additionally reverse the direction of every edge in the construction below.
\end{enumerate}
In the following, we assume that the class of all pure out-diamonds is $\oldC_1$.
To prove the statement, we first introduce the {\sc $k \times k$-Grid Tiling} problem, formally defined below.\\

\defparproblem{{\sc $k \times k$-Grid Tiling}}{Integers $k,n$ and $k^2$ nonempty set $S_{i,j} \subseteq [n] \times [n]$ where $i,j \in [k]$.}{$k$}{For each $i,j \in [k]$, does there exist an entry $(x_{i,j},y_{i,j}) \in S_{i,j}$ such that
\begin{itemize}
\item for every $i \in [k]$, $x_{i,1} = x_{i,2} = \ldots = x_{i,k}$ and
\item for every $j \in [k]$, $y_{1,j} = y_{2,j} = \ldots = y_{k,j}$?
\end{itemize}}

\bigskip

Throughout the paper, we assume that whenever given a set $S \subseteq [n] \times [n]$, $1 < x,y < n$ holds for every $(x,y) \in S$: it suffices to increase $n$ by two and replace $(x,y)$ by $(x+1,y+1)$ otherwise.

It was shown \cite[Theorem 14.28]{10.5555/2815661} that, under $\mathsf{ETH}$, {\sc $k \times k$-Grid Tiling} does not admit an algorithm running in time $f(k)\cdot n^{o(k)}$ for any computable function $f$. 
To prove \Cref{thm:diamond} for pure out-diamonds, 
we give a reduction which transforms an instance of {\sc $k \times k$-Grid Tiling} 
into an instance of (edge-weighted)\footnote{We argue at the end of the section that with polynomially bounded integer weights, the edge-weighted version of the problem reduces to its unweighted version.} {\sc Planar $\oldC_1$-Steiner Network} with $O(k^2)$ terminals.
To this end, we design three types of gadgets: 
the \emph{connector gadget}, the \emph{down main gadget} and the \emph{up main gadget}.
The reduction represents each set of the {\sc $k \times k$-Grid Tiling} instance with a copy of the up or down main gadget,
and uses the connector gadgets to further connect these gadgets (see \Cref{fig:reduction}).
The remainder of this section is organized as follows.
We first introduce the connector gadget and prove in \Cref{lem:cg} that it satisfies several desired properties.
We then introduce the down main gadget and prove in \Cref{lem:dmg} that it also satisfies several desired properties.
The up main gadget is introduced thereafter together with its symmetrical \Cref{lem:umg}.
We end the section with the precise description of the reduction and a proof of its correctness.\\

\noindent
\textbf{Connector gadget.} Given an integer $n > 0$, the connector gadget $CG_n$ is an edge-weighted planar digraph consisting of $n \times n$ grid, where the vertex lying at the intersection of column $i$ and row $j$ is denoted by $x_{i,j}$, and $2n+2$ additional vertices $p_1,\ldots,p_n,q_1,\ldots,q_n,p,q$. The adjacencies and edge weights are defined as follows (we fix $N = 3n$).
\begin{itemize}
\item \emph{Left source edges:} for every $j \in [n]$, there is an edge $(p_j,x_{1,j})$. Together these edges are called left source edges. The weight of each such edge is set to $N^2$.
\item \emph{Right source edges:} for every $j \in [n]$, there is an edge $(q_j,x_{n,j})$. Together these edges are called right source edges. The weight of each such edge is set to $N^2$.
\item \emph{Top sink edges:} for every $i \in [n]$, there is an edge $(x_{i,n},p)$. Together these edges are called top sink edges. The weight of each such edge is set to $N^2$.
\item \emph{Bottom sink edges:} for every $i \in [n]$, there is an edge $(x_{i,1},q)$. Together these edges are called bottom sink edges. The weight of each such edge is set to $N^2$.
\item The $n \times n$ grid is divided in two according to the diagonal $j = n+1-i$: the vertices on the top-right part (that is, the vertices $x_{i,j}$ such that $j+i \geq n+1$) are connected with $\leftarrow$ and $\uparrow$ edges, and the vertices of the bottom-left part (that is, the vertices $x_{i,j}$ such that $j+i \leq n+1$) are connected with $\rightarrow$ and $\downarrow$ edges. More specifically, we define the following edges.
\begin{itemize}
\item \emph{Up edges}: for every $i \in [n]$ and every $j \in [n-1]$ such that $j+i \geq n+1$, there is an edge $(x_{i,j},x_{i,j+1})$. Together these edges are called up edges.
\item \emph{Down edges:} for every $i \in [n]$ and every $j \in [n] \setminus \{1\}$ such that $j+i \leq n+1$, there is an edge $(x_{i,j},x_{i,j-1})$. Together these edges are called down edges.
\end{itemize}
We set the weight of every up/down edge to $N$.
\begin{itemize}
\item \emph{Left edges:} for every $i \in [n] \setminus \{1\}$ and every $j \in [n]$ such that $j+i > n+1$, there is an edge $(x_{i,j},x_{i-1,j})$. Together these edges are called left edges.
\item \emph{Right edges:} for every $i \in [n-1]$ and every $j \in [n]$ such that $j+i < n+1$, there is an edge $(x_{i,j},x_{i+1,j})$. Together these edges are called right edges.
\end{itemize}
We set the weight of every left/right edge to $1$.
\end{itemize}  
This concludes the construction of the connector gadget $CG_n$  (see \Cref{fig:connectorgad} for an illustration of the connector gadget for $n = 4$). In the following, we call the vertices $p_1,\ldots,p_n$ the \emph{left vertices}, the vertices $q_1,\ldots,q_n$ the \emph{right vertices} and the vertices $p,q$ the \emph{terminal vertices}. We further set 
\[
C_n^* = 4N^2+(n-1)N+n-1.
\]
A set $E \subseteq E(CG_n)$ satisfies the \emph{connectedness} property if the following hold in $E$:
\begin{itemize}
\item $p$ can be reached from some right vertex and from some left vertex;
\item $q$ can be reached from some right vertex and from some left vertex.
\end{itemize}
A set $E \subseteq E(CG_n)$ satisfying the connectedness property \emph{represents} an integer $j \in [n]$ if the only left source edge in $E$ is the one incident to $p_j$ and the only right source edge in $E$ is the one incident to $q_j$ (see \Cref{fig:connectorgad} for a set of edges representing 2).

\begin{figure}
\centering
\begin{tikzpicture}
\node[rectangle, fill=yellow, rounded corners,minimum width=5.3cm,minimum height=.25cm] at (3.5,2) {};
\node[rectangle, fill=yellow, rounded corners,minimum width=.25cm,minimum height=3.2cm] at (4,2.5) {};
\node[rectangle, fill=yellow, rounded corners,minimum width=.25cm,minimum height=1.35cm,rotate=27.5] at (3.75,4.5) {};
\node[rectangle, fill=yellow, rounded corners,minimum width=.25cm,minimum height=1.35cm,rotate=-27.5] at (3.75,.5) {};

\foreach \i in {2,...,5}
\foreach \j in {1,...,4}
{\node[circ] (\i\j) at (\i,\j) {};} 

\foreach \i in {2,...,4}
{\pgfmathtruncatemacro{\y}{6-\i}
  \foreach \j in {\y,...,2}
  {\pgfmathtruncatemacro{\z}{\j-1}
    \draw[->,>=stealth,thick] (\i\j) -- (\i\z);}}
    
\foreach \i in {5,...,3}
{\pgfmathtruncatemacro{\y}{6-\i}
  \foreach \j in {\y,...,3}
  {\pgfmathtruncatemacro{\z}{\j+1}
    \draw[->,>=stealth,thick] (\i\j) -- (\i\z);}}    
    
\foreach \j in {1,...,3}
{\pgfmathtruncatemacro{\y}{5-\j}
  \foreach \i in {2,...,\y}
  {\pgfmathtruncatemacro{\z}{\i+1}
   \draw[->,>=stealth,thick] (\i\j) -- (\z\j);}} 

\foreach \j in {4,...,2}
{\pgfmathtruncatemacro{\y}{7-\j}
  \foreach \i in {5,...,\y}
  {\pgfmathtruncatemacro{\z}{\i-1}
   \draw[->,>=stealth,thick] (\i\j) -- (\z\j);}} 

\foreach \j in {1,...,4}
{\node[circ,label=left:{\small $p_{\j}$}] (p\j) at (1,\j) {};}

\foreach \j in {1,...,4}
{\draw[->,>=stealth,thick] (p\j) -- (2\j);}

\foreach \j in {1,...,4}
{\node[circ,label=right:{\small $q_{\j}$}] (q\j) at (6,\j) {};}

\foreach \j in {1,...,4}
{\draw[->,>=stealth,thick] (q\j) -- (5\j);}

\node[circ,label=below:{\small $q$}] (q) at (3.5,0) {};

\foreach \i in {2,...,5}
{\draw[->,>=stealth,thick] (\i1) -- (q);}

\node[circ,label=above:{\small $p$}] (p) at (3.5,5) {};

\foreach \i in {2,...,5}
{\draw[->,>=stealth,thick] (\i4) -- (p);}
\end{tikzpicture}
\caption{The connector gadget for $n = 4$. A set of edges representing 2 is highlighted.}
\label{fig:connectorgad}
\end{figure}

\begin{lemma}
\label{lem:cg}
For any integer $n > 0$, the connector gadget $CG_n$ satisfies the following properties.
\begin{itemize}
\item[(1)] For every $j \in [n]$, there exists a set $E_j \subseteq E(CG_n)$ of weight $C_n^*$ representing $j$.
\item[(2)] If there exists a set $E \subseteq E(CG_n)$ of weight at most $C_n^*$ satisfying the connectedness property, then $E$ has weight exactly $C_n^*$ and represents some integer $j \in [n]$.
\end{itemize}
\end{lemma}

\begin{proof}
To prove (1), it suffices to take $E_j$ to be the union of the following sets of edges:
\begin{itemize}
\item $\{(p_j,x_{1,j}),(q_j,x_{n,j}),(x_{n+1-j,n},p),(x_{n+1-j,1},q)\}$; 
\item $\{(x_{i,j},x_{i+1,j})~|~1 \leq i \leq n-j\}$; 
\item $\{(x_{i,j},x_{i-1,j})~|~n+2-j \leq i \leq n\}$; 
\item $\{(x_{n+1-j,\ell}, x_{n+1-j,\ell+1})~|~j \leq \ell \leq n-1\}$; and
\item $\{(x_{n+1-j,\ell},x_{n+1-j,\ell -1})~|~2 \leq \ell \leq j\}$.
\end{itemize}
It is not difficult to see that $E_j$ represents $j$ and has weight $C_n^*$.\\

Next, suppose that $E \subseteq E(CG_n)$ is a set of weight at most $C_n^*$ satisfying the connectedness property. Let us show that $E$ has weight exactly $C_n^*$ and represents some integer $j \in [n]$. To this end, we prove the following claims.

\begin{claim}
\label{clm:sourcesinkedges}
$E$ contains exactly one left source edge, one right source edge, one top sink edge and one bottom sink edge.
\end{claim}

\begin{claimproof}
Since $E$ satisfies the connectedness property, it contains at least one left source edge, one right source edge, one top sink edge and one bottom sink edge. Now if $E$ contains at least two left source edges (the other cases are symmetric), then the weight of $E$ is at least $5N^2$; however, by definition,
\[
C_n^*  = 4N^2 + (n-1)N+ n-1 < 4N^2 + nN + nN < 5N^2
\]
as $N = 3n > 2n$, a contradiction.
\end{claimproof}

In the following, we let $p_{j_1}$ be the left vertex incident to the left source edge in $E$ and $q_{j_2}$ be the right vertex incident to the right source edge in $E$.

\begin{claim}
\label{clm:updownedges}
For every $j \in [n-1]$, $E$ contains exactly one (up or down) edge with one endvertex on row $j$ and one endvertex one row $j+1$. In particular, $E$ contains exactly $n-1$ up/down edges.
\end{claim}

\begin{claimproof}
First note that if there exists $j \in [n-1]$ such that $E$ contains no (up or down) edge with one endvertex on row $j$ and the other endvertex on row $j+1$, then either $j_1 \leq j$ in which case $p_{j_1}$ cannot reach $p$ in $E$, or $j_1 > j$ in which case $p_{j_1}$ cannot reach $q$ in $E$, a contradiction in both cases to the connectedness of $E$. Thus, for every $j \in [n-1]$, $E$ contains at least one (up or down) edge with one endvertex on row $j$ and the other endvertex on row $j+1$; in particular, $E$ contains at least $n-1$ up/down edges. Now if $E$ contains at least $n$ up/down edges then by \Cref{clm:sourcesinkedges}, the weight of $E$ is at least $4N^2 + nN$; however, by definition, 
\[
C_n^* = 4N^2 + (n-1)N + n-1 < 4N^2 + (n-1)N + n < 4N^2 + nN
\]
as $N = 3n > n$, a contradiction.
\end{claimproof}

\begin{claim}
\label{clm:leftrightedges}
$E$ contains exactly $n-1$ left/right edges.
\end{claim}

\begin{claimproof}
Observe first that for every $i \in [n-1]$, $E$ contains at least one (left or right) edge with one endvertex on column $i$ and the other endvertex on column $i+1$: indeed, if this were not the case for some $i \in [n-1]$, then every path in $CG_n[E]$ from $p_{j_1}$ to $p$ would be edge-disjoint from every path in $CG_n[E]$ from $q_{j_2}$ to $p$ and thus, $E$ would contain at least two top sink edges, a contradiction to \Cref{clm:sourcesinkedges}. Hence, for every $i \in [n-1]$, $E$ contains at least one (left or right) edge with one endvertex on column $i$ and the other endvertex on column $i+1$; in particular, $E$ contains at least $n-1$ left/right edges. Now if $E$ contains at least $n$ left/right edges then by Claims~\ref{clm:sourcesinkedges} and \ref{clm:updownedges}, the weight of $E$ is at least $4N^2 + (n-1)N + n > C_n^*$, a contradiction.
\end{claimproof}

Now by Claims~\ref{clm:sourcesinkedges}, \ref{clm:updownedges} and \ref{clm:leftrightedges}, $E$ has weight exactly $C_n^*$; we next show that $j_1 = j_2$ which would imply that $E$ represents $j_1$. Suppose to the contrary that $j_2 < j_1$ (the case where $j_2 > j_1$ is symmetric). By \Cref{clm:updownedges}, $E$ contains exactly one edge $e$ with one endvertex on row $j_2$ and one endvertex on row $j_2+1$; but then, either $e$ is an up edge in which case $p_{j_1}$ cannot reach $q$ in $E$, or $e$ is a down edge and $q_{j_2}$ cannot reach $p$ in $E$, a contradiction in both cases to the connectedness of $E$ which concludes the proof. 
\end{proof}

%=============================================================================================

\noindent
\textbf{Down Main Gadget.} Given an integer $n > 0$, the down main gadget $dMG_S$ represents a set $S \subseteq [n] \times [n]$\footnote{Recall that, by assumption, $1 < x,y< n$ holds for every $(x,y) \in S$.} and is constructed as follows. It is an edge-weighted planar digraph consisting of a $2n \times n^2$ grid, where the vertex lying at the intersection of column $i$ and row $j$ is denoted by $x_{i,j}$, and $6n$ additional vertices $\ell_1,\ldots,\ell_n,\ell'_1,\ldots,\ell'_n,r_1,\ldots,r_n,r'_1,\ldots,r'_n,t_1,\ldots,t_n$, $b_1,\ldots,b_n$. The adjacencies and edge weights are defined as follows (we fix $M=13n^2$). 

\begin{itemize}
\item \emph{Source edges:} for every $i \in [n]$, there is an edge $(t_i,x_{i,n^2})$. Together these edges are called source edges. The weight of each such edge is set to $M^5$.
\item \emph{Bottom sink edges:} for every $i \in [n]$, there is an edge $(x_{n+i,1},b_i)$. Together these edges are called bottom sink edges. The weight of each such edge is set to $M^5$.
\item \emph{Left sink edges:} for every $j \in [n]$, there is an edge $(\ell'_j,\ell_j)$ whose weight is set to $Mj$. Together these edges are called left sink edges.
\item  \emph{Right sink edges:} for every $j \in [n]$, there is an edge $(r'_j,r_j)$ whose weight is set to $M^2-Mj$. Together these edges are called right sink edges.
\item \emph{Left internal sink edges:} for every $j \in [n]$ and every $(j-1)n+1 \leq p \leq jn$, there is an edge $(x_{1,p},\ell'_j)$ whose weight is set to $p-(j-1)n$. For every fixed $j \in [n]$, these edges together are called the $j^{th}$ left internal sink edges.
\item \emph{Right internal sink edges:} for every $j \in [n]$ and every $(j-1)n + 1 \leq p \leq jn$, there is an edge $(x_{2n,p},r'_j)$ whose weight is set to $n+1-(p-(j-1)n)$. For every fixed $j \in [n]$, these edges together are called the $j^{th}$ right internal sink edges.
\item \emph{Right bridge edges:} for every $j \in [n]$ and every $(j-1)n + 1 \leq p \leq jn$, there is an edge $(x_{n,p},x_{n+1,p})$. For every fixed $j \in [n]$, these edges together are called the $j^{th}$ right bridge edges. The weight of each $j^{th}$ right bridge edge is set to $M^4$.
\item \emph{Downward bridge edges:} for every $j \in \{pn+1~|~p \in [n-1]\}$ and every $i \in [2n]$, there is an edge $(x_{i,j},x_{i,j-1})$. For every fixed $j \in \{pn+1~|~p \in [n-1]\}$, these edges together are called the $j^{th}$ downward bridge edges. The weight of each $j^{th}$ downward bridge edge is set to $M^3$.
\item \emph{Down edges:} for every $i \in [2n]$, every $j \in [n]$ and every $(j-1)n + 2 \leq p \leq jn$, there is an edge $(x_{i,p},x_{i,p-1})$. Together these edges are called down edges. The weight of each such edge is set to $4$.
\item \emph{Right edges:} for every $j \in [n]$, every $(j-1)n+1 \leq p \leq jn$ and every $ jn+1-p \leq q \leq 2n-1$ such that $q \neq n$, there is an edge $(x_{q,p},x_{q+1,p})$. Together these edges are called right edges. The weight of each such edge is set to $4$.
\item \emph{Left edges:} for every $j \in [n]$, every $(j-1)n+1 \leq p \leq jn-1$ and every $2 \leq q \leq jn+1-p$, there is an edge $(x_{q,p},x_{q-1,p})$. Together these edges are called left edges. The weight of each such edge is set to $4$.
\item \emph{Shortcut edges:} for every $s=(i,j) \in S$, we introduce two shortcut edges $e^\ell_s,e^r_s$ as follows. Set $p = jn+1$, then
\begin{itemize}
\item subdivide the edge $(x_{i,p+1-i},x_{i,p-i})$ by adding a vertex $y_{i,j}$ and the edge $(x_{i,p+1-i},y_{i,j})$ (of weight $3$) with the edge $(y_{i,j},x_{i,p-i})$ (of weight $1$);
\item subdivide the edge $(x_{n+i,p-i},x_{n+i,p-1-i})$ by adding a vertex $z_{i,j}$ and the edge $(x_{n+i,p-i},\allowbreak
z_{i,j})$ (of weight $3$) with the edge $(z_{i,j},x_{n+i,p-1-i})$ (of weight $1$). 
\end{itemize}
The introduced edges are called \emph{down subdivided edges}. Then
\begin{itemize}
\item $e^\ell_s= (y_{i,j},x_{i-1,p-i})$ and its weight is set to $2$;
\item $e^r_s = (z_{i,j},x_{n+i+1,p-i})$ and its weight is set to $2$.
\end{itemize}
The edges $e^\ell_s$ are called \emph{left} shortcut edges and the edges $e^r_s$ are called \emph{right} shortcut edges.  
\end{itemize}
This concludes the construction of the down main gadget $dMG_S$ (see \Cref{fig:downmaingad} for an illustration of the down main gadget with $n=4$ representing $S= \{(2,2),(2,3),(3,2)\}$). In the following, we call the vertices $\ell_1,\ldots,\ell_n$ the \emph{left vertices}, the vertices $r_1,\ldots,r_n$ the \emph{right vertices}, the vertices $t_1,\ldots,t_n$ the \emph{top vertices} and the vertices $b_1,\ldots,b_n$ the \emph{bottom vertices}. Furthermore, we set
\[
M_n^* = 2M^5 + M^4 + (n-1)M^3 + M^2 + (4n+1)(n+1) - 12.
\] 
We further let $V_\ell = \{x_{i,j}~|~1 \leq i \leq n \text{ and } 1 \leq j \leq n^2\} \cup \{y_{i,j}~|~(i,j) \in S\}$, $V_r = \{x_{i,j}~|~n+1 \leq i \leq 2n \text{ and } 1 \leq j \leq n^2\} \cup \{z_{i,j}~|~(i,j) \in S\}$ and for every $j \in [n]$, $V_j = \{x_{i,p}~|~1 \leq i \leq 2n \text{ and } (j-1)n + 1 \leq p \leq jn\}$. Note that, by construction, the following holds.

\begin{observation}
\label{obs:reach}
For every $j \in [n]$ and every $x_{i,p} \in V_j \cap V_\ell$, the following holds.
\begin{itemize}
\item If $n+1-i > p - (j-1)n$ then $x_{i,p}$ cannot reach a vertex $x_{i',p'} \in V_j$ such that $n+1-i' \leq p' - (j-1)n$.
\item If $n+1-i < p - (j-1)n$ then $x_{i,p}$ cannot reach a vertex $x_{i',p'} \in V_j$ such that $i' < i$ and $p' - (j-1)n > n+1-i$.
\end{itemize}
\end{observation}

\begin{figure}
\centering
\begin{tikzpicture}[scale=.8]
\node[rectangle, fill=yellow, rounded corners,minimum width=.25cm,minimum height=10.2cm] at (2,14.3) {};
\node[rectangle, fill=yellow, rounded corners,minimum width=4.25cm,minimum height=.25cm] at (4.5,8) {};
\node[rectangle, fill=yellow, rounded corners,minimum width=.25cm,minimum height=7.1cm] at (7,3.7) {};
\node[rectangle, fill=yellow, rounded corners,minimum width=1.5cm,minimum height=.25cm] at (-1.75,7.5) {};
\node[rectangle, fill=yellow, rounded corners,minimum width=1.5cm,minimum height=.25cm] at (11.75,7.5) {};
\node[rectangle, fill=yellow, rounded corners,minimum width=1.88cm,minimum height=.25cm,rotate=15] at (.05,7.75) {};
\node[rectangle, fill=yellow, rounded corners,minimum width=1.74cm,minimum height=.25cm,rotate=-15] at (10,7.75) {};
\node[rectangle, fill=yellow, rounded corners,minimum width=1cm,minimum height=.25cm] at (8.55,8) {};
\node[rectangle, fill=yellow, rounded corners,minimum width=1cm,minimum height=.25cm,rotate=32] at (1.5,8.3) {};
\node[rectangle, fill=yellow, rounded corners,minimum width=1cm,minimum height=.25cm,rotate=24] at (7.55,7.8) {};

\foreach \i in {1,...,4}
\foreach \j in {1,...,4}
{\node[circ] (\i\j) at (\i,\j) {};}

\foreach \i in {1,...,4}
\foreach \j in {6,...,9}
{\pgfmathtruncatemacro{\p}{\j-1}
\node[circ] (\i\p) at (\i,\j) {};}

\foreach \i in {1,...,4}
\foreach \j in {11,...,14}
{\pgfmathtruncatemacro{\p}{\j-2}
\node[circ] (\i\p) at (\i,\j) {};}

\foreach \i in {1,...,4}
\foreach \j in {16,...,19}
{\pgfmathtruncatemacro{\p}{\j-3}
\node[circ] (\i\p) at (\i,\j) {};}

\foreach \i in {6,...,9}
\foreach \j in {1,...,4}
{\pgfmathtruncatemacro{\p}{\i-1}
\node[circ] (\p\j) at (\i,\j) {};}

\foreach \i in {6,...,9}
\foreach \j in {6,...,9}
{\pgfmathtruncatemacro{\p}{\i-1}
\pgfmathtruncatemacro{\q}{\j-1}
\node[circ] (\p\q) at (\i,\j) {};}

\foreach \i in {6,...,9}
\foreach \j in {11,...,14}
{\pgfmathtruncatemacro{\p}{\i-1}
\pgfmathtruncatemacro{\q}{\j-2}
\node[circ] (\p\q) at (\i,\j) {};}

\foreach \i in {6,...,9}
\foreach \j in {16,...,19}
{\pgfmathtruncatemacro{\p}{\i-1}
\pgfmathtruncatemacro{\q}{\j-3}
\node[circ] (\p\q) at (\i,\j) {};}

\foreach \i in {6,...,9}
{\pgfmathtruncatemacro{\p}{\i-5}
\node[circ,label=below:{$b_\p$}] (b\p) at (\i,-.5) {};}

\foreach \i in {1,...,4}
{\node[circ,label=above:{$t_\i$}] (t\i) at (\i,20.5) {};}

\foreach \j in {0,...,3}
{\pgfmathsetmacro{\p}{\j*5+2.5}
\pgfmathtruncatemacro{\t}{\j+1}
\node[circ,label=below:{$\ell'_\t$}] (l'\t) at (-1,\p){};
\node[circ,label=left:{$\ell_\t$}] (l\t) at (-2.5,\p){};
\draw[->,>=stealth,thick] (l'\t) -- (l\t);}

\foreach \j in {0,...,3}
{\pgfmathsetmacro{\p}{\j*5+2.5}
\pgfmathtruncatemacro{\t}{\j+1}
\node[circ,label=below:{$r'_\t$}] (r'\t) at (11,\p){};
\node[circ,label=right:{$r_\t$}] (r\t) at (12.5,\p){};
\draw[->,>=stealth,thick] (r'\t) -- (r\t);}

\foreach \i in {1,...,8} 
\foreach \j in {2,...,16}
{\pgfmathtruncatemacro{\t}{\j-1}
\draw[->,>=stealth,thick] (\i\j) -- (\i\t);}

\foreach \j in {1,...,4}
{\pgfmathtruncatemacro{\a}{(\j-1)*4+1}
\pgfmathtruncatemacro{\b}{\j*4-1}
\foreach \p in {\a,...,\b}
{\pgfmathtruncatemacro{\d}{\p-(\j-1)*4}
\pgfmathtruncatemacro{\c}{5-\d}
\foreach \q in {2,...,\c}
{\pgfmathtruncatemacro{\t}{\q-1}
\draw[->,>=stealth,thick] (\q\p) -- (\t\p);}}}

\foreach \j in {1,...,4}
{\pgfmathtruncatemacro{\a}{(\j-1)*4+2}
\pgfmathtruncatemacro{\b}{\j*4}
\foreach \p in {\a,...,\b}
{\pgfmathtruncatemacro{\d}{\p-(\j-1)*4}
\pgfmathtruncatemacro{\c}{5-\d}
\foreach \q in {\c,...,3}
{\pgfmathtruncatemacro{\t}{\q+1}
\draw[->,>=stealth,thick] (\q\p) -- (\t\p);}}}

\foreach \j in {1,...,16}
\foreach \i in {4,...,7}
{\pgfmathtruncatemacro{\p}{\i+1}
\draw[->,>=stealth,thick] (\i\j) -- (\p\j);}

\foreach \i in {1,...,4}
{\pgfmathtruncatemacro{\t}{\i+4}
\draw[->,>=stealth,thick] (t\i) -- (\i16);
\draw[->,>=stealth,thick] (\t1) -- (b\i);}

\foreach \j in {1,...,4}
{\pgfmathtruncatemacro{\a}{(\j-1)*4+1}
\pgfmathtruncatemacro{\b}{\j*4}
\foreach \p in {\a,...,\b}
{\draw[->,>=stealth,thick] (1\p) -- (l'\j);
\draw[->,>=stealth,thick] (8\p) -- (r'\j);}}

\node[circr] (22l) at (2,8.6) {};
\draw[->,>=stealth,thick,red] (22l) --(17);
\node[circr] (22r) at (7,7.55) {};
\draw[->,>=stealth,thick,red] (22r) -- (77);

\node[circr] (23l) at (2,13.6) {};
\draw[->,>=stealth,thick,red] (23l) --(111);
\node[circr] (23r) at (7,12.55) {};
\draw[->,>=stealth,thick,red] (23r) -- (711);

\node[circr] (32l) at (3,7.6) {};
\draw[->,>=stealth,thick,red] (32l) --(26);
\node[circr] (32r) at (8,6.55) {};
\draw[->,>=stealth,thick,red] (32r) -- (86);
\end{tikzpicture}
\caption{The down main gadget $dMG_S$ with $n=4$ representing $S=\{(2,2),(2,3),(3,2)\}$ (the red edges are the shortcut edges). A set of edges representing $(2,2)$ is highlighted.}
\label{fig:downmaingad}
\end{figure}

\noindent
A set $E \subseteq E(dMG_S)$ satisfies the \emph{connectedness} property if the following hold in $E$: 
\begin{itemize}
\item a top vertex can reach a bottom vertex;
\item a top vertex can reach a right vertex; 
\item a top vertex can reach a left vertex.
\end{itemize}
A set $E \subseteq E(dMG_S)$ satisfying the connectedness property \emph{represents} a pair $(i,j) \in [n] \times [n]$ if the only source edge in $E$ is the one incident to $t_i$, the only bottom sink edge in $E$ is the one incident to $b_i$, the only left sink edge in $E$ is the one incident to $\ell_j$ and the only right sink edge in $E$ is the one incident to $r_j$ (see \Cref{fig:downmaingad} for a set of edges representing $(2,2)$).

\begin{lemma}
\label{lem:dmg}
For any $n > 0$ and any $S \subseteq [n] \times [n]$, the down main gadget $dMG_S$ satisfies the following properties.
\begin{itemize}
\item[(1)] For every $(i,j) \in S$, there exists a set $E_{i,j} \subseteq E(dMG_S)$ of weight $M_n^*$ representing $(i,j)$.
\item[(2)] If there exists a set $E \subseteq E(dMG_S)$ of weight at most $M_n^*$ satisfying the connectedness property, then $E$ has weight exactly $M_n^*$ and represents a pair $(i,j) \in S$.
\end{itemize}
\end{lemma}

\begin{proof}
To prove (1), it suffices to take $E_{i,j}$ to be the union of the following set of edges:
\begin{itemize}
\item $\{(t_i,x_{i,n^2}),(x_{n+i,1},b_i),(x_{0,jn+1-i},\ell'_j),(\ell'_j,\ell_j),(x_{2n+1,jn+1-i},r'_j),(r'_j,r_j)\}$;
\item $\{(x_{i,p},x_{i,p-1})~|~ jn+2-i\leq p \leq n^2\} \cup \{(x_{n+i,p},x_{n+i,p-1})~|~2 \leq p \leq jn+1-i\}$;
\item $\{(x_{p,jn+1-i},x_{p-1,jn+1-i})~|~ 1 \leq p \leq i-1\} \cup \{e^\ell_{(i,j)}\}$;
\item $\{(x_{p,jn+1-i},x_{p+1,jn+1-i})~|~i \leq p \leq 2n \text { and } p \neq n+i\} \cup \{e^r_{(i,j)}\}$.
\end{itemize}
It is not difficult to see that $E_{i,j}$ represents $(i,j)$ and has weight $M_n^*$.\\

Next, suppose that $E \subseteq E(dMG_S)$ is a set of weight at most $M_n^*$ satisfying the connectedness property. Let us show that $E$ has weight exactly $M_n^*$ and represents some $(i,j) \in S$. To this end, we first prove the following claims.

\begin{claim}
\label{clm:sourcesinkedges2}
$E$ contains exactly one source edge and one bottom sink edge.
\end{claim}

\begin{claimproof}
Since $E$ satisfies the connectedness property, it contains at least one source edge and one bottom sink edge. Now if $E$ contains at least two source edges, then the weight of $E$ is at least $3M^5$; however, by definition,
\begin{equation*}
\begin{split}
M_n^* &=  2M^5 + M^4 + (n-1)M^3 + M^2 + (4n+1)(n+1) - 12\\ 
& < 2M^5 + n^2M^4 + n^2M^4 + n^2M^4 + 9n^2M^4 < 3M^5
\end{split}
\end{equation*}
as $M = 13n^2 > 12n^2$, a contradiction. We conclude similarly if $E$ contains at least two bottom sink edges.
\end{claimproof}

In the following, we let $t_{i_1}$ be the top vertex incident to the source edge in $E$ and $b_{i_2}$ be the bottom vertex incident to the bottom sink edge in $E$. 

\begin{claim}
\label{clm:rightbridge}
$E$ contains exactly one right bridge edge.
\end{claim}

\begin{claimproof}
Since $t_{i_1}$ can reach $b_{i_2}$ in $E$, $E$ contains at least one right bridge edge. Now if $E$ contains at least two right bridge edges then by \Cref{clm:sourcesinkedges2}, the weight of $E$ is at least $2M^5 + 2M^4$; however, by definition,
\begin{equation*}
\begin{split}
M_n^* &= 2M^5 + M^4 + (n-1)M^3 + M^2 + (4n+1)(n+1) - 12\\
& < 2M^5 + M^4 + n^2M^4 + n^2M^4 + 9n^2M^4 \\
& < 2M^5 + 2M^4
\end{split}
\end{equation*}
as $M = 13n^2 > 11n^2$, a contradiction.
\end{claimproof}

\begin{claim}
\label{clm:downwardbridges}
For every $j \in [n-1]$, $E$ contains exactly one $j^{th}$ downward bridge edge. In particular, $E$ contains exactly $n-1$ downward bridge edges.
\end{claim}

\begin{claimproof}
Since $t_{i_1}$ can reach $b_{i_2}$ in $E$, $E$ contains at least one $j^{th}$ downward bridge edge for every $j \in [n-1]$. Now if $E$ contains at least $n$ downward bridge edges then by Claims~\ref{clm:sourcesinkedges2} and \ref{clm:rightbridge}, the weight of $E$ is at least $2M^5 + M^4 + nM^3$; however, by definition,
\begin{equation*}
\begin{split}
M_n^* &=  2M^5 + M^4 + (n-1)M^3 + M^2 + (4n+1)(n+1) - 12\\
& < 2M^5 + M^4 + (n-1)M^3 + n^2M^2 + 9n^2M^2 \\
& < 2M^5 + M^4 + nM^3
\end{split}
\end{equation*}
as $M = 13n^2 > 10n^2$, a contradiction.
\end{claimproof}

Since $E$ satisfies the connectedness property, it contains at least one left sink edge $(\ell'_{j_1},\ell_{j_1})$ and at least one right sink edge $(r'_{j_2},r_{j_2})$; we next show that $E$ contains in fact no other left or right sink edge.

\begin{claim}
\label{clm:j1=j2}
It holds that $j_1 = j_2$. Furthermore, $E$ contains a $j_1^{th}$ right bridge edge.
\end{claim}

\begin{claimproof}
Towards a contradiction, suppose that $j_1 \neq j_2$.
Assume first that $j_1 < j_2$ and let $P$ be a path from $t_{i_1}$ to $r_{j_2}$ in $dMG_S[E]$. Then $P$ must contain a $j^{th}$ right bridge edge $e$ for some $j \geq j_2$, and since $E$ contains exactly one right bridge edge by \Cref{clm:rightbridge}, it follows that every path in $dMG_S[E]$ from $t_{i_1}$ to $b_{i_2}$ contains $e$. But $t_{i_1}$ can reach $\ell_{j_1}$ in $E$ as well and so, $E$ must contain at least $n-j$ downward bridge edges for $t_{i_1}$ to reach the tail of $e$, plus $j -1$ downward bridge edges for the head of $e$ to then reach $b_{i_2}$, plus $j-j_1$ downward bridge edges for $t_{i_1}$ to also reach $\ell_{j_1}$, a contradiction to \Cref{clm:downwardbridges}. Suppose next that $j_2 < j_1$. Then since $E$ contains the edge $(\ell'_{j_1},\ell_{j_1})$ and the edge $(r'_{j_2},r_{j_2})$, it follows from Claims~\ref{clm:sourcesinkedges2}, \ref{clm:rightbridge} and \ref{clm:downwardbridges} that the weight of $E$ is at least $2M^5 + M^4 + (n-1)M^3 + Mj_1 + M^2 - Mj_2 \geq 2M^5 + M^4 + (n-1)M^3 + M^2 + M$. However, by definition,
\begin{equation*}
\begin{split}
M_n^* &= 2M^5 + M^4 + (n-1)M^3 + M^2 + (4n+1)(n+1) - 12\\
& <  2M^5 + M^4 + (n-1)M^3 + M^2 + 9n^2\\
& < 2M^5 + M^4 + (n-1)M^3 + M^2 + M
\end{split}
\end{equation*}
as $M = 13n^2 > 9n^2$, a contradiction. Thus, $j_1 = j_2$ and by arguing as above, we can show that $E$ must then contain a $j_1^{th}$ right bridge edge.
\end{claimproof}

In the following, we say that a path $P$ in $dMG_S$ \emph{crosses a row} $j$ if $V(P) \cap \{x_{i,j}~|~i \in [2n]\} \neq \emptyset$ and that $P$ \emph{crosses a column} $i$ if $V(P) \cap \{x_{i,j}~|~j \in [n^2]\} \neq \emptyset$. Furthermore, we define
\begin{itemize}
\item a \emph{$2d$-move} to be the union of two down subdivided edges $e,e'$ such that the head of $e$ coincides with the tail of $e'$;
\item a \emph{$d\ell$-move} to be the union of a down subdivided edge $e$ and a left shortcut edge $e'$ such that the head of $e$ coincides with the tail of $e'$;
\item a \emph{$dr$-move} to be the union of a down subdivided edge $e$ and a right shortcut edge $e'$ such that the head of $e$ coincides with the tail of $e'$. 
\end{itemize}
For convenience, we also call
\begin{itemize}
\item a right edge an \emph{$r$-move},
\item a left edge an \emph{$\ell$-move},
\item a down edge a \emph{$d$-move} and
\item a downward bridge edge a \emph{$Db$-move}.
\end{itemize}
Note that no $dr$-move is possible in $dMG_S[V_\ell]$ and that no $x$-move with $x \in \{d\ell,\ell\}$ is possible in $dMG_S[V_r]$. Furthermore, for any $x \in \{2d,r,\ell,d\}$, the weight of an $x$-move is 4 and for any $x \in \{dr,d\ell\}$, the weight of an $x$-move is 5. Given a path $P$ of $dMG_S$ and $x \in \{2d,d\ell,dr,r,\ell,d,Db\}$, we let $m(P,x)$ be the number of $x$-moves in $P$.\\

Let $(x_{1,j_0},\ell'_{j_1})$ be a $j_1^{th}$ left internal sink edge contained in $E$ (recall that $\ell_{j_1} \in V(E)$). Further let $x_{i_0,j_1n}$ be the head of the $j_1^{th}$ downward bridge in $E$ if $j_1 \leq n-1$ (recall that by \Cref{clm:downwardbridges}, $E$ contains exactly one such edge) and the head of the source edge in $E$ otherwise (note that in this case $i_0 = i_1$). Since $x_{i_0,j_1n}$ can reach $x_{1,j_0}$, the following holds.

\begin{observation}
\label{obs:i0j0}
$j_0 - (j_1-1)n \leq n+1-i_0$.
\end{observation}

Now by connectedness of $E$, there exists a path $P^t$ in $dMG_S[E]$ from the head of the source edge in $E$ to $x_{i_0,j_1n}$; we next lower bound the weight of $P^t$.

\begin{claim}
\label{clm:weightPt}
The weight of $P^t$ is at least $(n-j_1)M^3 + 4(n-1)(n-j_1) + |i_1 - i_0|$.
\end{claim}

\begin{claimproof}
Observe first that since by Claims~\ref{clm:rightbridge} and \ref{clm:j1=j2}, $E$ contains one $j_1^{th}$ right bridge edge and no other right bridge edge, necessarily $i_0 \leq n$. Now $P^t$ must cross every row $j_1n \leq j \leq n^2$ and thus, 
\[
m(P^t,Db) + m(P^t,d) + m(P^t,2d) + m(P^t,d\ell) = n^2 - j_1n.
\]
By \Cref{clm:downwardbridges}, $P^t$ contains exactly $n-j_1$ downward bridge edges, that is, $m(P^t,Db) = n - j_1$. We next distinguish cases depending on whether $i_1 \leq i_0$ or $i_1 > i_0$. Suppose first that $i_1 \leq i_0$. Then since $P^t$ must cross every column $i_1 \leq i \leq i_0$, $m(P^t,r) \geq i_0 - i_1$. It follows that the weight of $P^t$ is at least 
\begin{equation*}
\begin{split}
&~M^3m(P^t,Db) + 4m(P^t,d) + 4m(P^t,2d) + 5m(P^t,d\ell) + 4m(P^t,r) \\
& = (n-j_1)M^3 + 4m(P^t,d) + 4m(P^t,2d) + 5m(P^t,d\ell) + 4m(P^t,r) \\
& \geq (n-j_1)M^3 + 4(m(P^t,d) + m(P^t,2d) + m(P^t,d\ell)) + 4(i_0 - i_1) \\
& \geq (n-j_1)M^3 + 4(n^2 - j_1n - (n-j_1)) + 4(i_0 - i_1)\\
& \geq (n-j_1)M^3 + 4(n-1)(n-j_1) + (i_0 - i_1)
\end{split}
\end{equation*}
Second, suppose that $i_1 > i_0$. Then since $P^t$ must cross every column $i_0 \leq i \leq i_1$,
\[
m(P^t,\ell) + m(P^t,d\ell) = m(P^t,r) + i_1 - i_0
\]
and so, the weight of $P^t$ is at least
\begin{equation*}
\begin{split}
&~M^3m(P^t,Db) + 4m(P^t,d) + 4m(P^t,2d) + 5m(P^t,d\ell) + 4m(P^t,\ell) + 4m(P^t,r) \\
& = (n-j_1)M^3 + 4m(P^t,d) + 4m(P^t,2d) + 5m(P^t,d\ell) + 4m(P^t,\ell) + 4m(P^t,r) \\
& \geq (n-j_1)M^3 + 4(n^2 - j_1n - (n-j_1)) + m(P^t,d\ell) + 4m(P^t,\ell) + 4(m(P^t,\ell)\\
&+m(P^t,d\ell) - (i_1 - i_0))\\
& \geq (n-j_1)M^3 + 4(n-1)(n-j_1) + 5m(P^t,d\ell) + 8m(P^t,\ell) - 4(i_1 - i_0)\\
& \geq (n-j_1)M^3 + 4(n-1)(n-j_1) + 5(m(P^t,d\ell) + m(P^t,\ell)) - 4(i_1 - i_0)\\
& \geq (n-j_1)M^3 + 4(n-1)(n-j_1) + (i_1 - i_0)\\
\end{split}
\end{equation*}
as $m(P^t,d\ell) + m(P^t,\ell) \geq i_1 - i_0$ which proves our claim.
\end{claimproof}

Let $(x_{2n,j'_0},r'_{j_1})$ be a $j_1^{th}$ right internal sink edge contained in $E$ (recall that $r_{j_2} \in V(E)$ and $j_2 = j_1$ by \Cref{clm:j1=j2}). Further let $x_{i'_0,(j_1-1)n+1}$ be the tail of the $(j_1-1)^{th}$ downward bridge in $E$ if $j_1 >1$ (recall that by \Cref{clm:downwardbridges}, $E$ contains exactly one such edge) and the tail of the bottom sink edge contained in $E$ otherwise (note that in this case $i'_0 = n + i_2$). Then by connectedness of $E$, there exists a path $P^b$ in $dMG_S[E]$ from $x_{i'_0,(j_1-1)n+1}$ to the tail of the bottom sink edge in $E$; we next lower bound the weight of $P^b$.

\begin{claim}
\label{clm:weightPb}
The weight of $P^b$ is at least $(j_1-1)M^3 + 4(n -1)(j_1-1) + 4(n + i_2 - i'_0)$.
\end{claim}

\begin{claimproof}
Observe first that since by Claims~\ref{clm:rightbridge} and \ref{clm:j1=j2}, $E$ contains a $j_1^{th}$ right bridge edge and no other right bridge edge, necessarily $i'_0 \geq n+1$. Furthermore, since $dMG_S[V_r]$ contains no edge $r$-move, the following holds.

\begin{observation}
\label{obs:i'0}
$i'_0 \leq n+i_2$. 
\end{observation}

\noindent
Now $P^b$ must cross every column $i'_0 \leq i \leq n + i_2$ and every row $1 \leq j \leq (j_1-1)n+1$; thus 
\[
m(P^b,Db) + m(P^b,d) + m(P^b,2d) = (j_1-1)n \text{ and } m(P^b,r) + m(P^b,dr) = n + i_2 - i'_0.
\]
Since by \Cref{clm:downwardbridges}, $P^b$ contains exactly $j_1-1$ downward bridge edges, it follows that the weight of $P^b$ is at least 
\begin{equation*}
\begin{split}
&~M^3m(P^b,Db) + 4m(P^b,d) + 4m(P^b,2d) + 4m(P^b,r) + 5m(P^b,dr)\\
&= (j_1-1)M^3 + 4(m(P^b,d) + m(P^b,2d)) + 4m(P^b,r) + 5m(P^b,dr)\\
& \geq (j_1-1)M^3 + 4((j_1-1)n - (j_1-1)) + 4(m(P^b,r) + m(P^b,dr))\\
& \geq (j_1-1)M^3 + 4(j_1-1)(n - 1) + 4(n + i_2 - i'_0)
\end{split}
\end{equation*}
which proves our claim.
\end{claimproof}

Let $x_{n+1,t}$ be the head of the $j_1^{th}$ right bridge edge in $E$ (recall that by Claims~\ref{clm:rightbridge} and \ref{clm:j1=j2}, $E$ contains exactly one $j_1^{th}$ right bridge edge and no other right bridge edge). Then by connectedness of $E$, there exist in $dMG_S[E]$ a path $P$ and a path $P'$ from $x_{i_0,j_1n}$ to $x_{1,j_0}$ and from $x_{i_0,j_1n}$ to $x_{n,t}$ respectively; we next lower bound the weight of $P \cup P'$. To this end, denote by $x_{i',t'}$ the last vertex in $V(P)$ belonging to $V(P') \cap V_{j_1}$. 

\begin{observation}
\label{obs:i't'}
$\max\{t,j_0\} \leq t'$ and $n+1-i' \leq t' - (j_1-1)n$.
\end{observation}

Indeed, since $x_{i',t'}$ can reach both $x_{1,j_0}$ and $x_{n,t}$, $t' \geq j_0$ and $t' \geq t$. Now if $n+1-i' > t' - (j_1-1)n$ then by \Cref{obs:reach}, $x_{i',t'}$ cannot reach $x_{n,t}$ as $t - (j_1-1)n \geq 1 = n+1 -n$, a contradiction. 

\begin{observation}
\label{obs:i0i'}
$i_0 \leq i'$.
\end{observation}

Indeed, if $i_0 > i'$ then, in particular, $i_0 > 1$ and so, $n+1-i_0 < n = j_1n - (j_1-1)n$; but then, by \Cref{obs:reach}, $x_{i_0,j_1n}$ cannot reach $x_{i',t'}$ as $n+1-i_0 < n+1 - i' \leq t' - (j_1-1)n$ by \Cref{obs:i't'}, a contradiction. 

\begin{claim}
\label{clm:weightPP'}
The weight of $P \cup P'$ is at least
\[
4(n-1) + 4(t'-t) + 4(j_1n-j_0) + 4(i'-i_0) - 3(p_1 + m(P,d\ell))
\]
where 
$$
p_1 = \left\{
    \begin{array}{ll}
        1 & \mbox{if } (i',j_1) \in S, ~e^\ell_{i',j_1} \in E(P) \mbox{ and } x_{i',t'-1} \in V(P') \\
        0 & \mbox{otherwise.}
    \end{array}
\right.
$$
\end{claim}

\begin{claimproof}
Denote by $e$ the down or down subdivided edge with tail $x_{i',t'}$ (note that $e$ is a down subdivided edge if and only if $(i',j_1) \in S$). Then by definition of $x_{i',t'}$, the weight of $P \cup P'$ is at least the sum of the weights of $P'$ and $P[x_{i',t'},x_{1,j_0}]$ minus the weight of $e$ whenever $(i',j_1) \in S$ and $e \in E(P') \cap E(P[x_{i',t'},x_{1,j_0}])$. Let us therefore lower bound the weights of $P'$ and $P[x_{i',t'},x_{1,j_0}]$.

First note that since for any left shortcut edge $f$ with head $x_{p,q} \in V_j$, it holds that $n+1-p > q-(j-1)n$, $x_{p,q}$ cannot reach any vertex $x_{p',q'} \in V_j$ such that $n+1-p' \leq q' - (j-1)n$ by \Cref{obs:reach}. Further note that since the tail of $f$ can only be reached through $x_{p-1,q-1}$ where $n+1 - (p-1) < (q-1) - (j-1)n$, no vertex $x_{p',q'} \in V_j$ such that $n+1 - p' > q' - (j-1)n$ can reach the tail of $f$ by \Cref{obs:reach}. It follows that $P'$ and $P[x_{i_0,j_1n},x_{i',t'}]$ contain no left shortcut edge (recall that $n+1-i' \leq t' - (j_1-1)n$ by \Cref{obs:i't'}) and that $P[x_{i',t'},x_{1,j_0}]$ contains at most one left shortcut edge, that is, the following holds.

\begin{observation}
\label{obs:Pdl}
$m(P',d\ell) = m(P[x_{i_0,j_1n},x_{i',t'}], d\ell) = 0$ and $m(P,d\ell) \in \{0,1\}$. 
\end{observation}

\noindent
Now $P'$ must cross every row $t \leq j \leq j_1n$ and $P[x_{i',t'},x_{1,j_0}]$ must cross every row $j_0 \leq j \leq t'$ (recall that $\max\{t,j_0\} \leq t'$ by \Cref{obs:i't'}); thus
\begin{equation*}
\begin{split}
& m(P',d) + m(P',2d) = j_1n-t \text{ and } \\
& m(P[x_{i',t'}.x_{1,j_0}],d) + m(P[x_{i',t'}.x_{1,j_0}],2d) + m(P[x_{i',t'}.x_{1,j_0}],d\ell) = t'-j_0.
\end{split}
\end{equation*}
Similarly, $P'$ must cross every column $i_0 \leq i \leq n$ and $P[x_{i',t'},x_{1,j_0}]$ must cross every column $1 \leq i \leq i'$ (recall that $i_0 \leq i'$ by \Cref{obs:i0i'}); thus
\begin{equation*}
\begin{split}
& m(P',\ell) + m(P',r) \geq n-i_0 \text{ and }\\
& m(P[x_{i',t'}.x_{1,j_0}],\ell) + m(P[x_{i',t'}.x_{1,j_0}],d\ell) \geq i' - 1.
\end{split}
\end{equation*}
Now note that $p_1=1$ if and only if $(i',j_1) \in S$ and $e$ belongs to both $P$ and $P[x_{i',t'},x_{1,j_0}]$: indeed, if $p_1 =1$ then clearly $(i',j_1) \in S$ and $e \in E(P) \cap E(P[x_{i',t'},x_{1,j_0}])$. Conversely, if $(i',j_1) \in S$ and $e$ belongs to both $P$ and $P[x_{i',t'},x_{1,j_0}]$ then since by definition of $x_{i',t'}$, $x_{i',t'-1}$ cannot belong to both $P'$ and $P[x_{i',t'},x_{1,j_0}]$, and $P'$ contains no left shortcut edge by \Cref{obs:Pdl}, it must be that $P'$ contains $x_{i',t'-1}$ while $P[x_{i',t'},x_{1,j_0}]$ contains the left shortcut edge $e^\ell_{(i',j_1)}$, that is, $p_1=1$. It follows that the weight of $P \cup P'$ is at least
\begin{equation*}
\begin{split}
& 4(m(P',d) + m(P',2d)) + 4(m(P',\ell) + m(P',r)) + 4(m(P[x_{i',t'}.x_{1,j_0}],d) + m(P[x_{i',t'}.x_{1,j_0}],2d)) \\
+&~5(m(P[x_{i',t'}.x_{1,j_0}],d\ell) - p_1) + 2p_1 + 4m(P[x_{i',t'}.x_{1,j_0}],\ell)\\
\geq & ~4(j_1n-t) + 4(n-i_0) + 4(t'-j_0) + m(P[x_{i',t'}.x_{1,j_0}],d\ell) - 3p_1 \\
+&~4(i'-1 - m(P[x_{i',t'}.x_{1,j_0}],d\ell)) \\
\geq & ~4(n-1) + 4(t'-t) + 4(j_1n- j_0) + 4(i'-i_0) - 3(p_1+m(P[x_{i',t'}.x_{1,j_0}],d\ell))
\end{split}
\end{equation*}
which proves our claim as $m(P[x_{i',t'}.x_{1,j_0}],d\ell)=m(P,d\ell)$ by \Cref{obs:Pdl}.
\end{claimproof}

Note that if $p_1 = 1$ then, since $x_{i'-1,t'-1} \in V(P)$ can reach $x_{1,j_0}$ and $x_{i',t'-1} \in V(P')$ can reach $x_{n,t}$, the following holds.

\begin{observation}
\label{obs:p1}
If $p_1 = 1$ then $j_0 \leq t'-1$ and $t \leq t'-1$.
\end{observation} 

Now by connectedness of $E$, there exist in $dMG_S[E]$ a path $Q$ and a path $Q'$ from $x_{n+1,t}$ to $x_{i'_0,(j_1-1)n+1}$ and from $x_{n+1,t}$ to $x_{2n,j'_0}$, respectively; we next lower bound the weight of $Q \cup Q'$. To this end, denote by $x_{i'',t''}$ the last vertex in $V(Q')$ belonging to $V(Q) \cap V_{j_1}$. 

\begin{claim}
\label{clm:weightQQ'}
The weight of $Q \cup Q'$ is at least 
\[
4(n-2) + 4(i'_0-i'') + 4(t''-j'_0) + 4(t-(j_1-1)n) - 2p_2
\]
where 
$$
p_2 = \left\{
    \begin{array}{ll}
        1 & \mbox{if } (i''-n,j_1) \in S \mbox{ and } e^r_{i''-n,j_1} \in E(Q'),\\
        0 & \mbox{otherwise.}
    \end{array}
\right.
$$
\end{claim}

\begin{claimproof}
First note that since $x_{n+1,t}$ can reach $x_{i'',t''}$, and $x_{i'',t''}$ can reach both $x_{2n,j'_0}$ and $x_{i'_0,(j_1-1)n+1}$, the following holds.

\begin{observation}
\label{obs:t''}
$j'_0 \leq t'' \leq t$ and $i'' \leq i'_0$. 
\end{observation}

Now denote by $e$ the down or down subdivided edge with tail $x_{i'',t''}$ (note that $e$ is a down subdivided edge if and only if $(i''-n,j_1) \in S$). Then by definition of $x_{i'',t''}$, the weight of $Q \cup Q'$ is at least the sum of the weights of $Q$ and $Q'[x_{i'',t''},x_{2n,j'_0}]$ minus the weight of $e$ whenever $(i''-n,j_1) \in S$ and $e \in E(Q) \cap E(Q'[x_{i'',t''},x_{2n,j'_0}])$. Let us therefore lower bound the weights of $Q$ and $Q'[x_{i'',t''},x_{2n,j'_0}]$. 

Since $Q$ must cross every row $(j_1-1)n+1 \leq j \leq t$ and $Q'[x_{i'',t''},x_{2n,j'_0}]$ must cross every row $j'_0 \leq j \leq t''$ (recall that $j'_0 \leq t'' \leq t$ by \Cref{obs:t''}), it follows that
\begin{equation*}
\begin{split}
& m(Q,d) + m(Q,2d) = t-((j_1-1)n+1) \text{ and}\\
& m(Q'[x_{i'',t''},x_{2n,j'_0}],d) + m(Q'[x_{i'',t''},x_{2n,j'_0}],2d) = t'' - j'_0.
\end{split}
\end{equation*} 
Similarly, $Q$ must cross every column $n+1 \leq i \leq i'_0$ and $Q'[x_{i'',t''},x_{2n,j'_0}]$ must cross every column $i'' \leq i \leq 2n$ (recall that $i'' \leq i'_0$ by \Cref{obs:t''}); thus
\begin{equation*}
\begin{split}
& m(Q,r) + m(Q,dr) = i'_0 - (n+1) \text{ and}\\
& m(Q'[x_{i'',t''},x_{2n,j'_0}],r) + m(Q'[x_{i'',t''},x_{2n,j'_0}],dr) = 2n-i''.
\end{split}
\end{equation*}
Now note that $p_2 =1$ if and only if $(i''-n,j_1) \in S$ and $e$ belongs to both $Q$ and $Q'$: indeed, if $p_2=1$ then by definition $(i''-n,j_1) \in S$, $e \in E(Q')$ and $x_{i''+1,t''} \in V(Q')$ which implies, by definition of $x_{i'',t''}$, that $x_{i''+1,t''} \notin V(Q)$ and so, $e \in E(Q)$. Conversely, if $(i''-n,j_1) \in S$ and $e$ belongs to both $Q$ and $Q'$ then, by definition of $x_{i'',t''}$, it must be that $x_{i'',t''-1}$ belongs to one of $Q$ and $Q'$ while $x_{i''+1,t''}$ belongs to the other. However, if $x_{i''+1,t''}$ belongs to $Q$ (and so, $x_{i'',t''-1} \in V(Q')$) then, since $Q$ crosses row $(j_1-1)n+1$ and $Q'$ crosses column $2n$, necessarily $Q[x_{i''+1,t''},x_{i'_0,(j_1-1)n+1}] \cap Q'[x_{i'',t''-1},x_{2n,j'_0}] \cap V_{j_1} \neq \emptyset$, a contradiction to the definition of $x_{i'',t''}$. Thus, $x_{i''+1,t''} \in V(Q')$ and so, $e^r_{(i''-n,t'')} \in E(Q')$, that is, $p_2=1$. It follows that the weight of $Q \cup Q'$ is at least
\begin{equation*}
\begin{split}
&  4(m(Q,d) + m(Q,2d)) + 4m(Q,r) + 5m(Q,dr) + 4(m(Q'[x_{i'',t''},x_{2n,j'_0}],d) \\ 
+&~m(Q'[x_{i'',t''},x_{2n,j'_0}],2d))+4m(Q'[x_{i'',t''},x_{2n,j'_0}],r) + 5(m(Q'[x_{i'',t''},x_{2n,j'_0}],dr) - p_2) + 2p_2\\
\geq &~4(t-((j_1-1)n+1)) + 4(m(Q,r) + m(Q,dr)) + 4(t''-j'_0) + 4(2n-i'')\\ 
+&~ m(Q'[x_{i'',t''},x_{2n,j'_0}],dr) - 3p_2\\
\geq &~4(t-((j_1-1)n+1)) + 4(i'_0-(n+1)) + 4(t'' - j'_0) + 4(2n - i'')\\
+&~m(Q'[x_{i'',t''},x_{2n,j'_0}],dr) - 3p_2\\
\geq & ~4(n-2) + 4(i'_0-i'') + 4(t''-j'_0) + 4(t-(j_1-1)n) - 2p_2
\end{split}
\end{equation*}
which proves our claim.
\end{claimproof}

\noindent
In the following, we write $j_0 = (j_1-1)n+p_0$ and $j'_0 = (j_1-1)n+p'_0$ and further let 
\[
W = |i_1-i_0| + 4(n+i_2-i'') + 4(i'-i_0) + 4(t'-j_0) + 4(t''-j'_0) - (8 + 3(p_1 + m(P,d\ell)) + 2p_2).
\]
Observe that by Claims~\ref{clm:weightPt}, \ref{clm:weightPb}, \ref{clm:weightPP'} and \ref{clm:weightQQ'}, the weight of the union of $P^t,P^b,P,P',Q$ and $Q'$ is at least
\[
(n-1)M^3 + 4n(n+1) + W
\]
and so, by Claims~\ref{clm:sourcesinkedges2}, \ref{clm:rightbridge}, \ref{clm:downwardbridges} and \ref{clm:j1=j2}, the weight of $E$ is at least
\begin{equation}
\label{eq:weightE}
2M^5 + M^4 + (n-1)M^3 + M^2 + (4n+1)(n+1) + p_0 - p'_0 + W.
\end{equation}

\begin{claim}
\label{clm:i1=i2}
The following hold.
\begin{itemize}
\item[(1)] $i_1 = i_2$ and $(i_1,j_1) \in S$. 
\item[(2)] $W=-12$ and $p_0=p'_0$.
\end{itemize}
\end{claim}

\begin{claimproof}
Since the weight of $E$ is at most $M_n^*$, it follows from \Cref{eq:weightE} that
\begin{equation*}
\begin{split} 
M_n^* & = 2M^5 + M^4 + (n-1)M^3 + M^2 + (4n+1)(n+1) - 12\\
& \geq 2M^5 + M^4 + (n-1)M^3 + M^2 + (4n+1)(n+1) + p_0-p'_0 + W.
\end{split}
\end{equation*}
Thus, by definition of $W$,
\begin{equation}
\label{eq:upperbound}
\begin{split}
& p_0 - p'_0 + |i_1-i_0| + 4(n+i_2-i'') + 4(i'-i_0) + 4(t'-j_0) + 4(t''-j'_0)\\
&\leq -4 + 3(p_1 + m(P,d\ell)) + 2p_2 \\
& \leq 4
\end{split}
\end{equation}
\noindent
as $p_1,p_2 \in \{0,1\}$ by definition and $m(D,d\ell) \in \{0,1\}$ by \Cref{obs:Pdl}. Now suppose to the contrary that $p'_0 > p_0$ or, equivalently, that $j'_0 > j_0$. Then since $t' \geq t \geq t'' \geq j'_0$ by Observations~\ref{obs:i't'} and \ref{obs:t''}, $p'_0 - p_0 = j'_0 - j_0 \leq t' - j_0$ and so,
\begin{equation*}
\begin{split}
& p_0 - p'_0 + 4(t'-j_0) + |i_1-i_0| + 4(n+i_2-i'') + 4(i'-i_0)+ 4(t''-j'_0)\\
&\geq 3(t'-j_0) + |i_1-i_0| + 4(n+i_2-i'') + 4(i'-i_0) + 4(t''-j'_0).
\end{split}
\end{equation*} 
But $n+i_2 \geq i''$, $i' \geq i_0$ and $t' \geq t'' \geq j'_0 > j_0$ by Observations~\ref{obs:i'0}, \ref{obs:i't'}, \ref{obs:i0i'} and \ref{obs:t''}, hence
\begin{equation*}
3(t'-j_0) + |i_1-i_0| + 4(n+i_2-i'') + 4(i'-i_0) + 4(t''-j'_0) \geq 3.
\end{equation*} 
It then follows from \Cref{eq:upperbound} and the above that $p_1 = m(P,d\ell) = p_2=1$ and $t' = j_0 + 1$; however, by \Cref{obs:p1}, if $p_1 = 1$ then $t' \geq t+1$ and so, by \Cref{obs:i't'}, $j'_0 \leq t \leq t'-1 = j_0$, a contradiction. Thus, $p_0 \geq p'_0$ and so,
\begin{equation}
\label{eq:positif}
 p_0 - p'_0 + |i_1-i_0| + 4(n+i_2-i'') + 4(i'-i_0) + 4(t'-j_0) + 4(t''-j'_0) \geq 0
\end{equation}
\noindent 
by Observations~\ref{obs:i'0}, \ref{obs:i't'}, \ref{obs:i0i'}  and \ref{obs:t''}. Now suppose for a contradiction that $p_1=0$. Then by \Cref{eq:upperbound},
\[
 p_0 - p'_0 + |i_1-i_0| + 4(n+i_2-i'') + 4(i'-i_0) + 4(t'-j_0) + 4(t''-j'_0) \leq 1
\]
and so, by Observations~\ref{obs:i'0}, \ref{obs:i't'}, \ref{obs:i0i'} and \ref{obs:t''}, we must have $n+i_2 = i''$, $t' = j_0$, $t'' = j'_0$ and $i_0 = i'$. But then, by Observations~\ref{obs:i0j0} and \ref{obs:i't'}, $p_0 = n+1-i_0$ which implies that $P$ contains no left shortcut edge as $t'=j_0$; thus, by \Cref{eq:upperbound},  
\[
 p_0 - p'_0 + |i_1-i_0| + 4(n+i_2-i'') + 4(i'-i_0) + 4(t'-j_0) + 4(t''-j'_0) \leq -2,
\]
as $m(P,d\ell) = 0$, a contradiction to \Cref{eq:positif}. Hence, $p_1 = 1$ and so, $m(D,d\ell) = 1$ as well. Now by \Cref{obs:p1}, $t' \geq j_0+1$ and so, by \Cref{eq:upperbound},
\[
4 \leq 4(t'-j_0) \leq p_0 - p'_0 + |i_1-i_0| + 4(n+i_2-i'') + 4(i'-i_0)+ 4(t'-j_0) + 4(t''-j'_0) \leq 2 + 2p_2.
\]
Therefore, $p_2 = 1$ and $t'-1=j_0=j'_0=t''$, $i_1=i_0=i'$ and $n+i_2=i''$ by Observations~\ref{obs:i'0}, \ref{obs:i't'}, \ref{obs:i0i'} and \ref{obs:t''}; in particular, $W = -12$. Now since $p_1 = 1$, $(i',j_1) \in S$ by definition and so, $t' = j_1n+2 - i'$ by construction; and since $p_2 = 1$, $(i''-n,j_1) \in S$ by definition and so, $t''= j_1n+1-(i''-n)$ by construction. As $t' -1 = t''$, $i' = i_1$ and $i_2 = i''-n$, we conclude that $i_1 = i_2$ and $(i_1,j_1) \in S$. 
\end{claimproof}

It now follows from Claims~\ref{clm:j1=j2} and \ref{clm:i1=i2}(1) that $E$ represents the pair $(i_1,j_1) \in S$; and by combining \Cref{eq:weightE} with \Cref{clm:i1=i2}(2), we conclude that the weight of $E$ is at least $M_n^*$ and thus exactly $M_n^*$ which completes the proof.
\end{proof}

%=============================================================================================

\noindent
\textbf{Up Main Gadget.} Given an integer $n > 0$, the up main gadget $uMG_S$ represents a set $S \subseteq [n] \times [n]$\footnote{Recall that, by assumption, $1 < x,y< n$ holds for every $(x,y) \in S$.} and is constructed as follows. It is an edge-weighted planar digraph consisting of a $2n \times n^2$ grid, where the vertex lying at the intersection of column $i$ and row $j$ is denoted by $x_{i,j}$, and $6n$ additional vertices $\ell_1,\ldots,\ell_n,\ell'_1,\ldots,\ell'_n,r_1,\ldots,r_n,$ $r'_1,\ldots,r'_n,t_1,\ldots,t_n$, $b_1,\ldots,b_n$. The adjacencies and edge weights are defined as follows. 

\begin{itemize}
\item \emph{Source edges:} for every $i \in [n]$, there is an edge $(b_i,x_{n+i,1})$. Together these edges are called source edges. The weight of each such edge is set to $M^5$.
\item \emph{Top sink edges:} for every $i \in [n]$, there is an edge $(x_{i,n^2},t_i)$. Together these edges are called top sink edges. The weight of each such edge is set to $M^5$.
\item \emph{Left sink edges:} for every $j \in [n]$, there is an edge $(\ell'_j,\ell_j)$ whose weight is set to $Mj$. Together these edges are called left sink edges.
\item  \emph{Right sink edges:} for every $j \in [n]$, there is an edge $(r'_j,r_j)$ whose weight is set to $M^2-Mj$. Together these edges are called right sink edges. 
\item \emph{Left internal sink edges:} for every $j \in [n]$ and every $(j-1)n+1 \leq p \leq jn$, there is an edge $(x_{1,p},\ell''_j)$ whose weight is set to $p-(j-1)n$. For every fixed $j \in [n]$, these edges together are called the $j^{th}$ left internal sink edges.
\item \emph{Right internal sink edges:} for every $j \in [n]$ and every $(j-1)n + 1 \leq p \leq jn$, there is an edge $(x_{2n,p},r''_j)$ whose weight is set to $n+1-(p-(j-1)n)$. For every fixed $j \in [n]$, these edges together are called the $j^{th}$ right internal sink edges.
\item \emph{Left bridge edges:} for every $j \in [n]$, and every $(j-1)n + 1 \leq p \leq jn$, there is an edge $(x_{n+1,p},x_{n,p})$. For every fixed $j \in [n]$, these edges together are called the $j^{th}$ right bridge edges. The weight of each $j^{th}$ left bridge edge is set to $M^4$.
\item \emph{Upward bridge edges:} for every $j \in \{pn~|~p \in [n-1]\}$ and every $i \in [2n]$, there is an edge $(x_{i,j},x_{i,j+1})$. For every fixed $j \in \{pn~|~p \in [n-1]\}$, these edges together are called the $j^{th}$ upward bridge edges. The weight of each $j^{th}$ downward bridge edge is set to $M^3$.
\item \emph{Up edges:} for every $i \in [2n]$, every $j \in [n]$ and every $(j-1)n + 1 \leq p \leq jn-1$, there is an edge $(x_{i,p},x_{i,p+1})$. Together these edges are called up edges. The weight of each such edge is set to~$4$.
\item \emph{Left edges:} for every $j \in [n]$, every $(j-1)n+1 \leq p \leq jn-1$ and every $2 \leq q \leq jn+1-p$, there is an edge $(x_{q,p},x_{q-1,p})$; and for every $j \in [n]$, every $(j-1)n+1 \leq p \leq jn-1$ and every $2 \leq i \leq n$, there is an edge $(x_{i,p},x_{i-1,p}$. Together these edges are called left edges. The weight of each such edge is set to $4$.
\item \emph{Right edges:} for every $j \in [n]$, every $(j-1)n+2 \leq p \leq jn$ and every $jn+1-p \leq q \leq n-1$, there is an edge $(x_{n+q,p},x_{n+q+1,p})$. Together these edges are called right edges. The weight of each such edge is set to $4$.
\item \emph{Shortcut edges:} for every $s=(i,j) \in S$, we introduce two shortcut edges $e^\ell_s,e^r_s$ as follows. Set $p = jn+1$ then
\begin{itemize}
\item subdivide the edge $(x_{i,p-i},x_{i,p+1-i})$ by adding a vertex $y_{i,j}$ and the edge $(x_{i,p-i},y_{i,j})$ (of weight $3$) with the edge $(y_{i,j},x_{i,p+1-i})$ (of weight $1$);
\item subdivide the edge $(x_{n+i,p-1-i},x_{n+i,p-i})$ by adding a vertex $z_{i,j}$ and the edge 

$(x_{n+i,p-1-i},z_{i,j})$ (of weight $3$) with the edge $(z_{i,j},x_{n+i,p-i})$ (of weight $1$). 
\end{itemize}
The introduced edges are called \emph{down subdivided edges}. Then
\begin{itemize}
\item $e^\ell_s= (y_{i,j},x_{i-1,p-i})$ and its weight set to $2$;
\item $e^r_s = (z_{i,j},x_{n+i+1,p-i})$ and its weight set to $2$.
\end{itemize}
The edges $e^\ell_s$ are called \emph{left} shortcut edges and the edges $e^r_s$ are called \emph{right} shortcut edges.  
\end{itemize}
This concludes the construction of the up main gadget $uMG_S$ (see \Cref{fig:upmaingad} for an illustration of the up main gadget $uMG_S$ with $n=4$ representing $S= \{(2,2),(2,3),(3,2)\}$). We call the vertices $\ell_1,\ldots,\ell_n$ the \emph{left vertices}, the vertices $r_1,\ldots,r_n$ the \emph{right vertices}, the vertices $t_1,\ldots,t_n$ the \emph{top vertices} and the vertices $b_1,\ldots,b_n$ the \emph{bottom vertices}.\\

\begin{figure}
\centering
\begin{tikzpicture}[scale=.8]
\node[rectangle, fill=yellow, rounded corners,minimum width=.25cm,minimum height=10.2cm] at (2,14.3) {};
\node[rectangle, fill=yellow, rounded corners,minimum width=4.25cm,minimum height=.25cm] at (4.5,8) {};
\node[rectangle, fill=yellow, rounded corners,minimum width=.25cm,minimum height=7.1cm] at (7,3.7) {};
\node[rectangle, fill=yellow, rounded corners,minimum width=1.5cm,minimum height=.25cm] at (-1.75,7.5) {};
\node[rectangle, fill=yellow, rounded corners,minimum width=1.5cm,minimum height=.25cm] at (11.75,7.5) {};
\node[rectangle, fill=yellow, rounded corners,minimum width=1.88cm,minimum height=.25cm,rotate=15] at (.05,7.75) {};
\node[rectangle, fill=yellow, rounded corners,minimum width=1.74cm,minimum height=.25cm,rotate=-15] at (10,7.75) {};
\node[rectangle, fill=yellow, rounded corners,minimum width=1cm,minimum height=.25cm] at (8.55,8) {};
\node[rectangle, fill=yellow, rounded corners,minimum width=1cm,minimum height=.25cm,rotate=24] at (1.5,8.22) {};
\node[rectangle, fill=yellow, rounded corners,minimum width=1.15cm,minimum height=.25cm,rotate=29.4] at (7.55,7.73) {};

\foreach \i in {1,...,4}
\foreach \j in {1,...,4}
{\node[circ] (\i\j) at (\i,\j) {};}

\foreach \i in {1,...,4}
\foreach \j in {6,...,9}
{\pgfmathtruncatemacro{\p}{\j-1}
\node[circ] (\i\p) at (\i,\j) {};}

\foreach \i in {1,...,4}
\foreach \j in {11,...,14}
{\pgfmathtruncatemacro{\p}{\j-2}
\node[circ] (\i\p) at (\i,\j) {};}

\foreach \i in {1,...,4}
\foreach \j in {16,...,19}
{\pgfmathtruncatemacro{\p}{\j-3}
\node[circ] (\i\p) at (\i,\j) {};}

\foreach \i in {6,...,9}
\foreach \j in {1,...,4}
{\pgfmathtruncatemacro{\p}{\i-1}
\node[circ] (\p\j) at (\i,\j) {};}

\foreach \i in {6,...,9}
\foreach \j in {6,...,9}
{\pgfmathtruncatemacro{\p}{\i-1}
\pgfmathtruncatemacro{\q}{\j-1}
\node[circ] (\p\q) at (\i,\j) {};}

\foreach \i in {6,...,9}
\foreach \j in {11,...,14}
{\pgfmathtruncatemacro{\p}{\i-1}
\pgfmathtruncatemacro{\q}{\j-2}
\node[circ] (\p\q) at (\i,\j) {};}

\foreach \i in {6,...,9}
\foreach \j in {16,...,19}
{\pgfmathtruncatemacro{\p}{\i-1}
\pgfmathtruncatemacro{\q}{\j-3}
\node[circ] (\p\q) at (\i,\j) {};}

\foreach \i in {6,...,9}
{\pgfmathtruncatemacro{\p}{\i-5}
\node[circ,label=below:{$b_\p$}] (b\p) at (\i,-.5) {};}

\foreach \i in {1,...,4}
{\node[circ,label=above:{$t_\i$}] (t\i) at (\i,20.5) {};}

\foreach \j in {0,...,3}
{\pgfmathsetmacro{\p}{\j*5+2.5}
\pgfmathtruncatemacro{\t}{\j+1}
\node[circ,label=below:{$\ell'_\t$}] (l'\t) at (-1,\p){};
\node[circ,label=left:{$\ell_\t$}] (l\t) at (-2.5,\p){};
\draw[->,thick,>=stealth] (l'\t) -- (l\t);}

\foreach \j in {0,...,3}
{\pgfmathsetmacro{\p}{\j*5+2.5}
\pgfmathtruncatemacro{\t}{\j+1}
\node[circ,label=below:{$r'_\t$}] (r'\t) at (11,\p){};
\node[circ,label=right:{$r_\t$}] (r\t) at (12.5,\p){};
\draw[->,thick,>=stealth] (r'\t) -- (r\t);}

\foreach \i in {1,...,8} 
\foreach \j in {1,...,15}
{\pgfmathtruncatemacro{\t}{\j+1}
\draw[->,>=stealth,thick] (\i\j) -- (\i\t);}

\foreach \j in {1,...,4}
{\pgfmathtruncatemacro{\a}{(\j-1)*4+1}
\pgfmathtruncatemacro{\b}{\j*4-1}
\foreach \p in {\a,...,\b}
{\pgfmathtruncatemacro{\d}{\p-(\j-1)*4}
\pgfmathtruncatemacro{\c}{5-\d}
\foreach \q in {2,...,\c}
{\pgfmathtruncatemacro{\e}{\q+4}
\pgfmathtruncatemacro{\t}{\e-1}
\draw[->,>=stealth,thick] (\e\p) -- (\t\p);}}}

\foreach \j in {1,...,4}
{\pgfmathtruncatemacro{\a}{(\j-1)*4+2}
\pgfmathtruncatemacro{\b}{\j*4}
\foreach \p in {\a,...,\b}
{\pgfmathtruncatemacro{\d}{\p-(\j-1)*4}
\pgfmathtruncatemacro{\c}{5-\d}
\foreach \q in {\c,...,3}
{\pgfmathtruncatemacro{\d}{\q+4}
\pgfmathtruncatemacro{\t}{\d+1}
\draw[->,>=stealth,thick] (\d\p) -- (\t\p);}}}

\foreach \j in {1,...,16}
\foreach \i in {1,...,4}
{\pgfmathtruncatemacro{\p}{\i+1}
\draw[<-,>=stealth,thick] (\i\j) -- (\p\j);}

\foreach \i in {1,...,4}
{\pgfmathtruncatemacro{\t}{\i+4}
\draw[<-,>=stealth,thick] (t\i) -- (\i16);
\draw[<-,>=stealth,thick] (\t1) -- (b\i);}

\foreach \j in {1,...,4}
{\pgfmathtruncatemacro{\a}{(\j-1)*4+1}
\pgfmathtruncatemacro{\b}{\j*4}
\foreach \p in {\a,...,\b}
{\draw[->,>=stealth,thick] (1\p) -- (l'\j);
\draw[->,>=stealth,thick] (8\p) -- (r'\j);}}

\node[circr] (22l) at (2,8.45) {};
\draw[->,>=stealth,thick,red] (22l) --(17);
\node[circr] (22r) at (7,7.4) {};
\draw[->,>=stealth,thick,red] (22r) -- (77);

\node[circr] (23l) at (2,13.45) {};
\draw[->,>=stealth,thick,red] (23l) --(111);
\node[circr] (23r) at (7,12.4) {};
\draw[->,>=stealth,thick,red] (23r) -- (711);

\node[circr] (32l) at (3,7.45) {};
\draw[->,>=stealth,thick,red] (32l) --(26);
\node[circr] (32r) at (8,6.4) {};
\draw[->,>=stealth,thick,red] (32r) -- (86);
\end{tikzpicture}
\caption{The up main gadget $uMG_S$ with $n=4$ representing $S=\{(2,2),(2,3),(3,2)\}$ (the red edges are the shortcut edges). A set of edges representing $(2,2)$ is highlighted.}
\label{fig:upmaingad}
\end{figure}

\noindent
A set $E \subseteq E(uMG_S)$ satisfies the \emph{connectedness} property if the following hold in $E$: 
\begin{itemize}
\item a bottom vertex can reach a top vertex;
\item a bottom vertex can reach a left vertex;
\item a bottom vertex can reach a right vertex.
\end{itemize}
A set $E \subseteq E(uMG_S)$ satisfying the connectedness property \emph{represents} a pair $(i,j) \in [n] \times [n]$ if the only source edge in $E$ is the one incident to $b_i$, the only top sink edge in $E$ is the one incident to $t_i$, the only left sink edge in $E$ is the one incident to $\ell_j$ and the only right sink edge in $E$ is the one incident to $r_j$ (see \Cref{fig:upmaingad} for a set of edges representing $(2,2)$). Symmetrical to \Cref{lem:dmg}, we have the following.

\begin{lemma}
\label{lem:umg}
For any $n > 0$ and any $S \subseteq [n] \times [n]$, the up main gadget $uMG_S$ satisfies the following properties.
\begin{itemize}
\item[(1)] For every $(i,j) \in S$, there exists a set $E_{i,j} \subseteq E(uMG_S)$ of weight $M_n^*$ representing $(i,j)$.
\item[(2)] If there exists a set $E \subseteq E(uMG_S)$ of weight at most $M_n^*$ satisfying the connectedness property then $E$ has weight exactly $M_n^*$ and represents a pair $(i,j) \in S$.
\end{itemize}
\end{lemma}

%=============================================================================================

\noindent
\textbf{Reduction.} Given an instance $(k,n,\{S_{i,j}~|~i,j \in [k]\})$ of \textsc{Grid Tiling}, we construct an equivalent instance $(G,T,D)$ of (edge-weighted) \textsc{Planar $\oldC_1$-Steiner Network} where $G$ is defined as follows (see \Cref{fig:reduction}). 
\begin{itemize}
\item We introduce a total of $k^2$ (down/up) main gadgets and $k(k+1)$ connector gadgets.
\item For every set $S_{i,j}$ of the \textsc{Grid Tiling} instance such that $i$ is odd, we introduce an up main gadget $uMG_{i,j}$ representing $S_{i,j}$. The up main gadget $uMG_{i,j}$ is surrounded by two connector gadgets: $CG_{i,j}$ lying to its left and $CG_{i+1,j}$ lying to its right. We identify each right vertex of $CG_{i,j}$ with the left vertex of $uMG_{i,j}$ of the same index, and each left vertex of $CG_{i+1,j}$ with the right vertex of $uMG_{i,j}$ of the same index. Furthermore, for every $j \in [k-1]$, $uMG_{i,j}$ lies above $uMG_{i,j+1}$ and we identify each bottom vertex of $uMG_{i,j}$ with the top vertex of $uMG_{i,j+1}$ of the same index.
\item For every set $S_{i,j}$ of the \textsc{Grid Tiling} instance such that $i$ is even, we introduce a down main gadget $dMG_{i,j}$ representing $S_{i,j}$. The down main gadget $dMG_{i,j}$ is surrounded by two connector gadgets: $CG_{i,j}$ lying to its left and $CG_{i+1,j}$ lying to its right. We identify each right vertex of $CG_{i,j}$ with the left vertex of $dMG_{i,j}$ of the same index, and each left vertex of $CG_{i+1,j}$ with the right vertex of $dMG_{i,j}$ of the same index. Furthermore, for every $j \in [k-1]$, $dMG_{i,j}$ lies above $dMG_{i,j+1}$ and we identify each bottom vertex of $dMG_{i,j}$ with the top vertex of $dMG_{i,j+1}$ of the same index.
\item We introduce $k$ terminals $t_1,\ldots,t_k$ and add the following edges of weight 0: for every odd $i \in [k]$, we add an edge from each top vertex of $uMG_{i,1}$ to $t_i$ and for every even $i \in [k]$, we add an edge from each bottom vertex of $dMG_{i,k}$ to $t_i$.
\item We introduce two vertices $r_1$ and $r_2$ and add the following edges of weight 0: there is an edge from $r_1$ to every terminal $t_i$ such that $i \in [k]$ is odd; for every even $i \in [k]$, there is an edge from $r_1$ to every top vertex of $dMG_{i,1}$; for every $j \in [k]$, there is an edge from $r_1$ to every left vertex of $CG_{1,j}$; and if $k$ is odd then for every $j \in [k]$, there is an edge from $r_1$ to every right vertex of $CG_{k+1,j}$. Similarly, there is an edge from $r_2$ to every terminal $t_i$ such that $i$ is even; for every odd $i \in [k]$, there is an edge from $r_1$ to every bottom vertex of $dMG_{i,k}$; and if $k$ is even then for every $j \in [k]$, there is an edge from $r_1$ to every right vertex of $CG_{k+1,j}$.
\end{itemize} 
This concludes the construction of $G$. The set $T$ of terminals consists of the union of the two terminal vertices in each connector gadget and $\{r_1,r_2\} \cup \{t_i~|~i \in [k]\}$ (note that $|T| = 2k(k+1) + k + 2$). The demand graph $D$ is the pure out-diamond on vertex set $T$ where $r_1$ and $r_2$ are the two vertices of in-degree 0. In the following, we let 
\[
W_n^* = k^2M_n^* + k(k+1)C_n^*.   
\]

\begin{figure}
\centering
\begin{tikzpicture}
\foreach \i in {0,...,3}
\foreach \j in {0,...,2}
{\pgfmathsetmacro{\a}{\i*4}
\pgfmathsetmacro{\b}{\a+.5}
\pgfmathsetmacro{\c}{\j*3}
\pgfmathsetmacro{\d}{\c+1.5}
\filldraw[fill=lightgray,very thick] (\a,\c) rectangle (\b,\d);
\pgfmathsetmacro{\e}{\a+.25}
\pgfmathsetmacro{\f}{\c-.35}
\pgfmathtruncatemacro{\p}{\i+1}
\pgfmathtruncatemacro{\q}{\j+1}
\node[draw=none] at (\e,\f) {$CG_{\p,\q}$};
\pgfmathsetmacro{\f}{\c+.2}
\node[circ] at (\e,\f) {};
\pgfmathsetmacro{\g}{\c+1.3}
\node[circ] at (\e,\g) {};
}

\foreach \j in {0,...,2}
{\pgfmathsetmacro{\c}{\j*3}
\pgfmathsetmacro{\d}{\c+1.5}
\filldraw[color=red!60,fill=red!5,very thick] (1.5,\c) rectangle (3,\d);
\pgfmathsetmacro{\e}{\c+.75}
\pgfmathtruncatemacro{\q}{\j+1}
\node[draw=none] at (2.25,\e) {$uMG_{1,\q}$};
\pgfmathsetmacro{\f}{\c+.3}
\draw[->,thick,>=stealth] (1.4,\f) -- (.6,\f);
\pgfmathsetmacro{\g}{\c+.78}
\node[draw=none,rotate=90] at (1,\g) {$\cdots$};
\pgfmathsetmacro{\h}{\c+1.2}
\draw[->,thick,>=stealth] (1.4,\h) -- (.6,\h);
\draw[->,thick,>=stealth] (3.1,\f) -- (3.9,\f);
\node[draw=none,rotate=90] at (3.5,\g) {$\cdots$};
\draw[->,thick,>=stealth] (3.1,\h) -- (3.9,\h);
}

\foreach \i in {0,2}
\foreach \j in {0,...,1}
{\pgfmathsetmacro{\a}{\i*4}
\pgfmathsetmacro{\b}{\a+1.8}
\pgfmathsetmacro{\c}{\j*3+1.6}
\pgfmathsetmacro{\d}{\c+1.3}
\draw[->,thick,>=stealth] (\b,\c) -- (\b,\d);
\pgfmathsetmacro{\e}{\b+.9}
\draw[->,thick,>=stealth] (\e,\c) -- (\e,\d);
\pgfmathsetmacro{\f}{\b+.45}
\pgfmathsetmacro{\g}{\c+.65}
\node[draw=none] at (\f,\g) {$\cdots$};}

\foreach \i in {1}
\foreach \j in {0,...,1}
{\pgfmathsetmacro{\a}{\i*4}
\pgfmathsetmacro{\b}{\a+1.8}
\pgfmathsetmacro{\c}{\j*3+1.6}
\pgfmathsetmacro{\d}{\c+1.3}
\draw[<-,thick,>=stealth] (\b,\c) -- (\b,\d);
\pgfmathsetmacro{\e}{\b+.9}
\draw[<-,thick,>=stealth] (\e,\c) -- (\e,\d);
\pgfmathsetmacro{\f}{\b+.45}
\pgfmathsetmacro{\g}{\c+.65}
\node[draw=none] at (\f,\g) {$\cdots$};}

\foreach \j in {0,...,2}
{\pgfmathsetmacro{\c}{\j*3}
\pgfmathsetmacro{\d}{\c+1.5}
\filldraw[color=blue!60,fill=blue!5,very thick] (5.5,\c) rectangle (7,\d);
\pgfmathsetmacro{\e}{\c+.75}
\pgfmathtruncatemacro{\q}{\j+1}
\node[draw=none] at (6.25,\e) {$dMG_{2,\q}$};
\pgfmathsetmacro{\f}{\c+.3}
\draw[->,thick,>=stealth] (5.4,\f) -- (4.6,\f);
\pgfmathsetmacro{\g}{\c+.78}
\node[draw=none,rotate=90] at (5,\g) {$\cdots$};
\pgfmathsetmacro{\h}{\c+1.2}
\draw[->,thick,>=stealth] (5.4,\h) -- (4.6,\h);
\draw[->,thick,>=stealth] (7.1,\f) -- (7.9,\f);
\node[draw=none,rotate=90] at (7.5,\g) {$\cdots$};
\draw[->,thick,>=stealth] (7.1,\h) -- (7.9,\h);}

\foreach \j in {0,...,2}
{\pgfmathsetmacro{\c}{\j*3}
\pgfmathsetmacro{\d}{\c+1.5}
\filldraw[color=red!60,fill=red!5,very thick] (9.5,\c) rectangle (11,\d);
\pgfmathsetmacro{\e}{\c+.75}
\pgfmathtruncatemacro{\q}{\j+1}
\node[draw=none] at (10.25,\e) {$uMG_{3,\q}$};
\pgfmathsetmacro{\f}{\c+.3}
\draw[->,thick,>=stealth] (9.4,\f) -- (8.6,\f);
\pgfmathsetmacro{\g}{\c+.78}
\node[draw=none,rotate=90] at (9,\g) {$\cdots$};
\pgfmathsetmacro{\h}{\c+1.2}
\draw[->,thick,>=stealth] (9.4,\h) -- (8.6,\h);
\draw[->,thick,>=stealth] (11.1,\f) -- (11.9,\f);
\node[draw=none,rotate=90] at (11.5,\g) {$\cdots$};
\draw[->,thick,>=stealth] (11.1,\h) -- (11.9,\h);}

\node[circ,label=left:{$t_1$}] (t1) at (2.25,8.5) {};
\draw[->,thick,>=stealth] (1.8,7.6) -- (t1); 
\draw[->,thick,>=stealth] (2.7,7.6) -- (t1); 
\node[draw=none] at (2.25,7.8) {$\cdots$};
\node[circ,label=left:{$t_2$}] (t2) at (6.25,-1) {};
\draw[->,thick,>=stealth] (5.8,-.1) -- (t2); 
\draw[->,thick,>=stealth] (6.7,-.1) -- (t2); 
\node[draw=none] at (6.25,-.3) {$\cdots$};
\node[circ,label=right:{$t_3$}] (t3) at (10.25,8.5) {};
\draw[->,thick,>=stealth] (9.8,7.6) -- (t3); 
\draw[->,thick,>=stealth] (10.7,7.6) -- (t3); 
\node[draw=none] at (10.25,7.8) {$\cdots$};

\node[circ,label=above:{$r_1$}] (r1) at (6.25,9.5) {};
\draw[->,thick,>=stealth] (r1) -- (t1);
\draw[->,thick,>=stealth] (r1) -- (t3);
\draw[<-,thick,>=stealth] (5.8,7.6) -- (6.25,8.5); 
\draw[<-,thick,>=stealth] (6.7,7.6) -- (6.25,8.5); 
\node[draw=none] at (6.25,7.8) {$\cdots$};
\draw[thick] (r1) -- (6.25,8.5);

\draw[->,thick,>=stealth] (-.5,6.75) -- (-.1,7.2);
\draw[->,thick,>=stealth] (-.5,6.75) -- (-.1,6.3);
\node[draw=none,rotate=90] at (-.2,6.77) {$\cdots$};
\draw[thick] (-.5,8.85) -- (-.5,6.75);
\draw[thick] (13,8.85) -- (13,6.75);

\draw[->,thick,>=stealth] (-.75,3.75) -- (-.1,4.2);
\draw[->,thick,>=stealth] (-.75,3.75) -- (-.1,3.3);
\node[draw=none,rotate=90] at (-.2,3.77) {$\cdots$};
\draw[thick] (-.75,8.825) -- (-.75,3.75);
\draw[thick] (13.25,8.85) -- (13.25,3.75);

\draw[->,thick,>=stealth] (-1,.75) -- (-.1,1.2);
\draw[->,thick,>=stealth] (-1,.75) -- (-.1,.3);
\node[draw=none,rotate=90] at (-.2,.77) {$\cdots$};
\draw[thick] (r1) -- (-1,8.8) -- (-1,.75);
\draw[thick] (r1) -- (13.5,8.8) -- (13.5,.75);

\draw[->,thick,>=stealth] (13,6.75) -- (12.6,7.2);
\draw[->,thick,>=stealth] (13,6.75) -- (12.6,6.3);
\node[draw=none,rotate=90] at (12.7,6.77) {$\cdots$};

\draw[->,thick,>=stealth] (13.25,3.75) -- (12.6,4.2);
\draw[->,thick,>=stealth] (13.25,3.75) -- (12.6,3.3);
\node[draw=none,rotate=90] at (12.7,3.77) {$\cdots$};

\draw[->,thick,>=stealth] (13.5,.75) -- (12.6,1.2);
\draw[->,thick,>=stealth] (13.5,.75) -- (12.6,.3);
\node[draw=none,rotate=90] at (12.7,.77) {$\cdots$};

\node[circ,label=below:{$r_2$}] (r2) at (6.25,-2) {};
\draw[->,thick,>=stealth] (r2) -- (t2);
\draw[<-,thick,>=stealth] (1.8,-.1) -- (2.25,-1); 
\draw[<-,thick,>=stealth] (2.7,-.1) -- (2.25,-1);
\node[draw=none] at (2.25,-.3) {$\cdots$};
\draw[thick] (r2) -- (2.25,-1); 
\draw[<-,thick,>=stealth] (9.8,-.1) -- (10.25,-1); 
\draw[<-,thick,>=stealth] (10.7,-.1) -- (10.25,-1);
\node[draw=none] at (10.25,-.3) {$\cdots$};
\draw[thick] (r2) -- (10.25,-1); 
\end{tikzpicture}
\caption{An illustration of the reduction from \textsc{Grid Tiling} to  \textsc{Planar $\oldC_1$-Steiner Network} with $k=3$ (the black vertices are the terminals).}
\label{fig:reduction}
\end{figure}

\begin{lemma}
The \textsc{Grid Tiling} instance $(k,n,\{S_{i,j}~|~i,j \in [k]\})$ has a solution if and only if the (edge-weighted) \textsc{Planar $\oldC_1$-Steiner Network} instance $(G,T,D)$ has a solution of weight at most $W_n^*$.
\end{lemma}

\begin{proof}
Assume first that the instance $(k,n,\{S_{i,j}~|~i,j \in [k]\})$ of \textsc{Grid Tiling} has a solution, that is, for every $i,j \in [k]$, there is an entry $(x_{i,j},y_{i,j}) \in S_{i,j}$ such that
\begin{itemize}
\item for every $i \in [k]$, $x_{i,1} = x_{i,2} = \ldots = x_{i,k} = \alpha_i$ and
\item for every $j \in [k]$, $y_{1,j} = y_{2,j} = \ldots = y_{k,j} = \beta_j$.
\end{itemize}
We construct a solution $E$ for the \textsc{Planar $\oldC_1$-Steiner Network} instance $(G,T,D)$ of weight at most $W_n^*$ as follows. Include in $E$
\begin{itemize}
\item for every odd $i \in [k]$, the edge $(r_1,t_i)$ of weight 0 and for every even $i \in [k]$, the edge $(r_2,t_i)$ of weight 0;
\item for every even $i \in [k]$, the edge from $r_1$ to the top vertex of index $\alpha_i$ in $dMG_{i,1}$ of weight 0 and for every odd $i \in [k]$, the edge from $r_2$ to the bottom vertex of index $\alpha_i$ in $dMG_{i,k}$ of weight 0;  
\item for every $j \in [k]$, the edge from $r_1$ to the left vertex of index $\beta_j$ in $CG_{1,j}$ of weight 0;
\item if $k$ is odd then for every $j \in [k]$, the edge from $r_1$ to the right vertex of index $\beta_j$ in $CG_{k+1,j}$ of weight 0 and if $k$ is even then for every $j \in [k]$, the edge from $r_2$ to the right vertex of $\beta_j$ in $CG_{k+1,j}$ of weight 0;
\item for every $j \in [k]$ and every $i \in [k]$, the set $E^C_{i,j} \subseteq E(CG_{i,j})$ of weight $C_n^*$ representing $\beta_j$ whose existence is guaranteed by \Cref{lem:cg}(1);
\item for every $j \in [k]$ and every odd $i \in [k]$, the set $E^M_{i,j} \subseteq E(uMG_{i,j})$ of weight $M_n^*$ representing $(\alpha_i,\beta_j)$ whose existence is guaranteed by \Cref{lem:umg}(1); and
\item for every $j \in [k]$ and every even $i \in [k]$, the set $E^M_{i,j} \subseteq E(dMG_{i,j})$ of weight $M_n^*$ representing $(\alpha_i,\beta_j)$ whose existence is guaranteed by \Cref{lem:dmg}(1).
\end{itemize}
It is not difficult to see that the weight of $E$ is $k^2M_n^* + k(k+1)C_n^* = W_n^*$ and that by the connectedness of the sets $E^C_{p,j}$ and $E^M_{i,j}$ for every $i,j \in [k]$ and $p \in [k+1]$, $r_1$ and $r_2$ can reach every terminal in $T \setminus \{r_1,r_2\}$ in $E$.\\ 

Conversely, assume that the \textsc{Planar $\oldC_1$-Steiner Network} instance $(G,T,D)$ has a solution $E$ of weight at most $W_n^*$. We contend that for every $i \in [k+1]$ and every $j \in [k]$, $E \cap E(CG_{i,j})$ satisfies the connectedness property. Indeed, if this is not the case for some $i \in [k+1]$ and $j\in [k]$, say no left vertex of $V(CG_{i,j})$ can reach the terminal $p \in V(CG_{i,j})$ in $E$ (the other cases are symmetric) then either $i$ is even in which case $r_2$ cannot reach $p$ in $E$, or $i$ is odd in which case $r_1$ cannot reach $p$ in $E$, a contradiction in both cases. Similarly, the restriction of $E$ to any (down/up) main gadget satisfies the connectedness property: indeed, if for some $j \in [k]$ and for some even $i \in [k]$ (we argue similarly if $i$ is odd), $E \cap E(dMG_{i,j})$ does not satisfy the connectedness property, then either no top vertex of $dMG_{i,j}$ can reach a bottom vertex of $dMG_{i,j}$ in $E$ in which case $r_2$ cannot reach $t_i$ in $E$; or no top vertex of $dMG_{i,j}$ can reach a left (or right) vertex of $dMG_{i,j}$ in $E$ in which case $r_2$ (or $r_1$) cannot reach the terminal vertices of $CG_{i,j}$ in $E$, a contradiction in both cases.

Next, we argue that the weight of the restriction of $E$ to any connector gadget is $C_n^*$ and that the weight of the restriction of $E$ to any (down/up) main gadget is $M_n^*$. To this end, let $c$ and $C$ be the number of connector gadgets whose weight in $E$ is at most $C_n^*$ and greater than $C_n^*$, respectively. Then $c+C=k(k+1)$ and by \Cref{lem:cg}(2), any connector gadget whose weight in $E$ is at most $C_n^*$ has in fact a weight of exactly $C_n^*$ in $E$. Similarly, let $m$ and $M$ be the number of (down/up) main gadgets whose weight in $E$ is at most $M_n^*$ and greater than $M_n^*$ respectively. Then $m + M = k^2$ and by Lemmas~\ref{lem:dmg} and \ref{lem:umg}, any (down/up) main gadget whose weight in $E$ is at most $M_n^*$ has in fact a weight of exactly $M_n^*$ in $E$. Now by definition of $W_n^*$, 
\begin{equation*}
\begin{split}
W_n^*& = k^2M_n^* + k(k+1)C_n^*\\
& \geq mM_n^* + M(M_n^*+1) + cC_n^* + C(C_n^*+1)\\
& = k^2M_n^* + M + k(k+1)C_n^* + C
\end{split}
\end{equation*}  
which implies that $M = C = 0$. Thus, every connector gadget has weight $C_n^*$ in $E$ and every (down/up) main gadget has weight $M_n^*$ in $E$. From Lemmas~\ref{lem:cg}(2), \ref{lem:dmg}(2) and \ref{lem:umg}(2), it then follows that
\begin{itemize}
\item for every $j \in [k]$ and every $i \in [k+1]$, the restriction of $E$ to the connector gadget $CG_{i,j}$ represents an integer $\beta'_{i,j} \in [n]$;
\item for every $j \in [k]$ and every even $i \in [k]$, the restriction of $E$ to the down main gadget $dMG_{i,j}$ represents a pair $(\alpha_{i,j},\beta_{i,j}) \in [n] \times [n]$; and
\item for every $j \in [k]$ and every odd $i \in [k]$, the restriction of $E$ to the up main gadget $uMG_{i,j}$ represents a pair $(\alpha_{i,j},\beta_{i,j}) \in [n] \times [n]$.
\end{itemize}
Let us show that for every $i,j \in [k]$ the entries $(\alpha_{i,j},\beta_{i,j}) \in S_{i,j}$ form a solution to the \textsc{Grid Tiling} instance $(k,n,\{S_{i,j}~|~i,j \in [k]\})$ which if true, would conclude the proof. To this end, we first prove that for every $i,j \in [k]$, $\beta'_{i,j} = \beta_{i,j}$. Consider an even $i \in [k]$. Then by \Cref{lem:dmg}(2), the only left sink edge in $E$ incident to a left vertex of $dMG_{i,j}$ is the one incident to the left vertex of index $\beta_{i,j}$; and by \Cref{lem:cg}(2), the only right source edge in $E$ incident to a right vertex of $CG_{i,j}$ is the one incident to the right vertex of index $\beta'_{i,j}$. Thus, if $\beta'_{i,j} \neq \beta_{i,j}$ then $r_1$ cannot reach the terminal vertices of $CG_{i,j}$ in $E$, a contradiction. We conclude similarly if $i$ is odd. 

Second, we show that for every $i,j \in [k]$, $\beta'_{i+1,j} = \beta_{i,j}$. Consider an even $i \in [k]$. Then by \Cref{lem:dmg}(2), the only right sink edge in $E$ incident to a right vertex of $dMG_{i,j}$ is the one incident to the right vertex of index $\beta_{i,j}$; and by \Cref{lem:cg}(2), the only left source edge in $E$ incident to a left vertex of $CG_{i+1,j}$ is the one incident to the left vertex of index $\beta'_{i,j}$. Thus, if $\beta'_{i,j} \neq \beta_{i,j}$ then $r_2$ cannot reach the terminal vertices of $CG_{i+1,j}$ in $E$, a contradiction. We conclude similarly if $i$ is odd. 

It follows from the above that for every $i \in [k-1]$ and every $j \in [k]$, $\beta_{i+1,j} = \beta'_{i+1,j} = \beta_{i,j}$; we next show that for every $i \in [k]$ and every $j \in [k-1]$, $\alpha_{i,j} = \alpha_{i,j+1}$. Consider an even $i \in [k]$. Then by \Cref{lem:dmg}(2), the only bottom sink edge in $E$ incident to a bottom vertex of $dMG_{i,j}$ is the one incident to the bottom vertex of index $\alpha_{i,j}$; and by \Cref{lem:dmg}(2), the only source edge in $E$ incident to a top vertex of $dMG_{i,j+1}$ is the one incident to the top vertex of index $\alpha_{i,j+1}$. Thus, if $\alpha_{i,j} \neq \alpha_{i,j+1}$ then $r_1$ cannot reach the terminal $t_i$ in $E$, a contradiction. We conclude similarly if $i$ is odd. Therefore, the entries $(\alpha_{i,j},\beta_{i,j}) \in S_{i,j}$ form a solution to the \textsc{Grid Tiling} instance $(k,n,\{S_{i,j}~|~i,j \in [k]\})$ as claimed.
\end{proof}

Let us finally explain how to get rid of the edge-weights (we use the same trick as in \cite{ChitnisFHM20}).
We replace every edge $(x,y)$ of weight $w$ in the instance $(G,T,D)$ of (edge-weighted) {\sc Planar $\mathcal{C}_1$-Steiner Network} constructed above,
with a directed path from $x$ to $y$ of length $w\cdot n + 1$ where $n = |V(G)|$. 
We let $G'$ be the resulting graph.
Then similarly to \cite[Theorem A.1]{ChitnisFHM20}, we can show that the instance $(G,T,D)$ of edge-weighted {\sc Planar $\mathcal{C}_1$-Steiner Network} has a solution of weight at most $W$ if and only if the instance $(G',T,D)$ of {\sc Planar $\mathcal{C}_1$-Steiner Network} has a solution of size at most $Wn + n$.
